\makeatletter
%% ===== 「名前と操作」の型（tools/kata9-int.md、組 V9-2）。頭文字 kixg は章ごとに置きかえる =====
\colorlet{kixgNc}{bkOrange!85!black}% 名前の枠の線（点線）
\colorlet{kixgSc}{bkTheme}% 実際の操作の枠の線（実線）
\colorlet{kixgSb}{bkTheme!8}% 実際の操作の地
\tikzset{
  kixg dan/.style={font=\sffamily\bfseries\footnotesize, rounded corners=1.2mm, inner xsep=1.6mm, inner ysep=0.5mm, line width=0.9pt},
  kixg dan1/.style={kixg dan, draw=kixgNc, dash pattern=on 2.2pt off 1.4pt, fill=white, text=bkInk},
  kixg dan2/.style={kixg dan, draw=kixgSc, fill=kixgSb, text=bkInk},
  kixg dan3/.style={kixg dan, draw=bkTheme, fill=bkTheme, text=bkOn},
  kixg name/.style={draw=kixgNc, line width=0.9pt, dash pattern=on 2.2pt off 1.4pt, rounded corners=1.2mm, inner xsep=1.2mm, inner ysep=1.2mm, fill=white},
  kixg seki/.style={draw=kixgSc, line width=0.9pt, rounded corners=1.2mm, inner xsep=1.2mm, inner ysep=1.2mm, fill=kixgSb},
  kixg tag/.style={font=\sffamily\bfseries\footnotesize, text=bkInk, inner sep=0.4mm},
  kixg q/.style={diamond, aspect=4.6, draw=bkYellow!75!black, line width=0.9pt, fill=bkYellowBg, font=\sffamily\footnotesize, align=flush center, inner sep=0.2mm, text=bkInk},
  kixg do/.style={rectangle, draw=\bkTT, line width=0.9pt, fill=white, rounded corners=0.6mm, font=\sffamily\footnotesize, align=flush center, inner xsep=1.5mm, inner ysep=1.2mm, text=bkInk},
  kixg go/.style={-{Stealth[length=2.2mm]}, line width=0.9pt, bkInk!75, shorten <=1.5mm, shorten >=1.5mm},
  kixg yn/.style={font=\sffamily\bfseries\footnotesize, text=bkInk, fill=white, inner sep=0.5mm},
}
% 実際の操作の4種類の名前（実線の札の字）  1 ふつうの定積分　2 表を引く　3 部分分数　4 行列の計算
\def\kixg@op#1{\ifcase#1\or ふつうの定積分\or 表を引く\or 部分分数\or 行列の計算\fi}
% 段の札（例題の中で）  \kixgdan{1|2|3}
\newcommand\kixgdan[1]{\tikz[baseline=(d.base)]\node[kixg dan#1] (d) {\vphantom{あ}\ifcase#1\or 名前\or 直す\or 楽をする\fi};}
% 直す段の札に、どの道具かをそえる（例題の中で）  \kixgnaosu{1|2|3|4} →「直す：表を引く」
\newcommand\kixgnaosu[1]{\tikz[baseline=(d.base)]\node[kixg dan2] (d) {\vphantom{あ}直す：\kixg@op{#1}};}
% 実線の札だけ（表の「手順」の頭・本文）  \kixgop{1|2|3|4}
\newcommand\kixgop[1]{\tikz[baseline=(d.base)]\node[kixg dan2] (d) {\vphantom{あ}\kixg@op{#1}};}
% 段の見出し（本文で）  \kixgdanhead{1|2|3}{見出し}
\newcommand\kixgdanhead[2]{\par\addvspace{8pt}\Needspace{7\baselineskip}\noindent
  \tikz[baseline=(d.base)]\node[kixg dan#1, font=\sffamily\bfseries\small, inner ysep=0.8mm] (d) {\vphantom{あ}\ifcase#1\or ①\,名前の段\or ②\,直す段\or ③\,楽をする段\fi};%
  \hspace{0.7\zw}{\sffamily\bfseries\color{\bkTT}#2}\par\nobreak\vspace{3pt}\noindent\ignorespaces}
% 場合の札  \kixgcase{C}（′ は \kixgcase{A$'$}）
\newcommand\kixgcase[1]{\tikz[baseline=(c.base)]\node[fill=bkTheme, text=bkOn, font=\sffamily\bfseries\small, rounded corners=0.8mm, inner xsep=1mm, inner ysep=0.6mm, minimum width=5.4mm] (c) {\vphantom{Ag}#1};}
% 例題の1行め  \kixgbaai{ことば}{9-2}{C}（場合が2つなら {D,B}：通った順）
\newcommand\kixgbaai[3]{{\sffamily\bfseries 場合：}#1（表#2の\kixg@cases#3,\@nil）}
\def\kixg@cases#1,#2\@nil{\kixgcase{#1}\ifx\relax#2\relax\else と\kixg@cases#2\@nil\fi}
% 式の中の「名前」（点線）と「実際の操作」（実線）
\newcommand\kixgN[1]{\tikz[baseline=(n.base)]\node[kixg name] (n) {$\displaystyle #1$};}
\newcommand\kixgNt[1]{\tikz[baseline=(n.base)]{\node[kixg name] (n) {$\displaystyle #1$};\node[kixg tag, anchor=north] at ([yshift=-0.8mm]n.south) {名前};}}
\newcommand\kixgS[1]{\tikz[baseline=(n.base)]\node[kixg seki] (n) {$\displaystyle #1$};}
% 札つきの実線の枠  \kixgSt{1|2|3|4}{式}（札の字：ふつうの定積分・表を引く・部分分数・行列の計算）
\newcommand\kixgSt[2]{\tikz[baseline=(n.base)]{\node[kixg seki] (n) {$\displaystyle #2$};\node[kixg tag, anchor=north] at ([yshift=-0.8mm]n.south) {\kixg@op{#1}};}}
% 場合分けの表の題の帯  \kixgcasetitle{9-2}{題}
\newcommand\kixgcasetitle[2]{\par\addvspace{6pt}\Needspace{8\baselineskip}\noindent\begin{tikzpicture}[baseline=(t.base)]
  \node[fill=bkTheme, text=bkOn, font=\sffamily\bfseries\small, rounded corners=1mm, inner xsep=2mm, inner ysep=1mm] (t) {場合分けの表\,#1};
  \node[anchor=west, font=\sffamily\bfseries, text=\bkTT, inner sep=0pt] at ([xshift=2mm]t.east) {#2};
  \end{tikzpicture}\par\nobreak\vspace{3pt}}
% 名前と操作の対応表の小さな絵  \kixgpic{1〜7}
%  1 フーリエ係数（1周期で cos をかける）　2 複素係数（回る矢印）　3 フーリエ変換（−∞〜∞）
%  4 ラプラス変換（0から、e^{-st} で小さく）　5 逆ラプラス変換（s の平面の極）　6 伝達関数（箱）　7 e^{At}（行列とうずまき）
\newcommand\kixgpic[1]{\begin{tikzpicture}[x=1mm, y=1mm, baseline=(current bounding box.center)]
  \ifcase#1\or
  \draw[bkGray, line width=0.4pt] (0,4) -- (14,4);
  \draw[bkInk, line width=0.8pt] (1,4) -- (1,7) -- (7,7) -- (7,1) -- (13,1) -- (13,4);
  \draw[bkTheme, line width=0.8pt, dash pattern=on 1.6pt off 0.9pt, domain=1:13, samples=40, smooth] plot (\x,{4+2.6*cos(30*(\x-1))});
  \draw[bkInk, line width=0.5pt] (1,-0.6) -- (1,-1.6) -- (13,-1.6) -- (13,-0.6);
  \or
  \draw[bkGray, line width=0.4pt] (2,4) -- (12,4) (7,-0.5) -- (7,8.5);
  \draw[bkInk, line width=0.7pt] (7,4) circle (3.6);
  \draw[bkTheme, line width=1.1pt, -{Stealth[length=1.6mm]}] (7,4) -- ({7+3.6*cos(-40)},{4+3.6*sin(-40)});
  \draw[bkInk, line width=0.6pt, -{Stealth[length=1.3mm]}] ({7+5*cos(10)},{4+5*sin(10)}) arc (10:-60:5);
  \or
  \draw[bkInk, line width=0.6pt, {Stealth[length=1.4mm]}-{Stealth[length=1.4mm]}] (0,1) -- (14,1);
  \draw[bkInk, line width=0.8pt] (1.5,1) -- (5.5,1) -- (5.5,6.5) -- (8.5,6.5) -- (8.5,1);
  \draw[bkTheme, line width=0.8pt, dash pattern=on 1.6pt off 0.9pt, domain=1.5:12.5, samples=50, smooth] plot (\x,{1+2.2*(1+cos(80*(\x-7)))/2});
  \or
  \draw[bkInk, line width=0.6pt, -{Stealth[length=1.4mm]}] (1,1) -- (14,1);
  \draw[bkInk, line width=0.6pt] (1,0) -- (1,8.5);
  \fill[bkTheme!25, domain=1:13, samples=40, smooth] (1,1) -- plot (\x,{1+6.5*exp(-0.3*(\x-1))}) -- (13,1) -- cycle;
  \draw[bkTheme, line width=0.8pt, domain=1:13, samples=40, smooth] plot (\x,{1+6.5*exp(-0.3*(\x-1))});
  \or
  \draw[bkGray, line width=0.4pt] (0,4) -- (14,4) (10,-0.5) -- (10,8.5);
  \foreach \kixgx/\kixgy in {3/4,6.5/6.8,6.5/1.2} {\draw[bkRed, line width=1pt] ({\kixgx-1},{\kixgy-1}) -- ({\kixgx+1},{\kixgy+1}) ({\kixgx-1},{\kixgy+1}) -- ({\kixgx+1},{\kixgy-1});}
  \or
  \draw[bkInk, line width=0.6pt, -{Stealth[length=1.4mm]}] (0,4) -- (4,4);
  \draw[bkTheme, line width=0.8pt, fill=bkTheme!12] (4.4,1) rectangle (9.6,7);
  \node[font=\footnotesize, text=bkInk, inner sep=0pt] at (7,4) {$G$};
  \draw[bkInk, line width=0.6pt, -{Stealth[length=1.4mm]}] (10,4) -- (14,4);
  \or
  \draw[bkInk, line width=0.6pt] (1.6,0.5) -- (0.6,0.5) -- (0.6,7.5) -- (1.6,7.5) (5.4,0.5) -- (6.4,0.5) -- (6.4,7.5) -- (5.4,7.5);
  \foreach \kixgx/\kixgy in {2.3/2.3,4.7/2.3,2.3/5.7,4.7/5.7} {\fill[bkInk] (\kixgx,\kixgy) circle (0.5);}
  \draw[bkTheme, line width=0.8pt, domain=0:600, samples=80, smooth] plot ({10.5+3.4*exp(-0.0035*\x)*cos(\x)},{4+3.4*exp(-0.0035*\x)*sin(\x)});
  \fi
\end{tikzpicture}}
\makeatother
%% ===== 表計算・2つの世界の図の型（tools/kata9.md、組 V9-1）。頭文字 kixg =====
\makeatletter
\colorlet{kixgN}{bkOrange!16}% 1行め（成分の番号・周波数）のマス
\colorlet{kixgK}{bkGreen!16}% 2行め（成分の大きさ）のマス
\colorlet{kixgS}{bkYellow!45}% 合計・注目するマス
\newcommand\kixgfx[2]{\tikz[baseline=(f.base)]{\node[draw=bkLine, line width=0.6pt, fill=white, minimum height=5.5mm, inner xsep=2mm, font=\sffamily\footnotesize, minimum width=12mm] (c) {#1};
  \node[anchor=west, draw=bkLine, line width=0.6pt, fill=white, minimum height=5.5mm, inner xsep=2mm, minimum width=78mm, text width=74mm, align=left, font=\ttfamily\footnotesize] (f) at ([xshift=1mm]c.east) {#2};}}
\colorlet{kixgW}{bkBlue!80!black}% 時間の世界の線
\colorlet{kixgF}{bkOrange!85!black}% 周波数の世界の棒
\colorlet{kixgP}{bkRed!80!black}% s の世界の点（×）
\pgfplotsset{
  kixg base/.style={bkplot, scale only axis, width=47mm, height=28mm, axis lines=left, grid=none,
    every axis x label/.style={at={(axis description cs:0.5,0)}, anchor=north, yshift=-5mm, font=\sffamily\small, text=bkInk},
    every axis y label/.style={at={(axis description cs:0,1)}, anchor=south, yshift=1.5mm, font=\sffamily\small, text=bkInk}},
  kixg wave/.style={kixg base, extra y ticks={0}, extra y tick labels={0}, extra y tick style={grid=major, major grid style={bkInk!45, line width=0.5pt}}},
  kixg spec/.style={kixg base, ymin=0},
  kixg splane/.style={bkplot, scale only axis, width=47mm, height=47mm, axis equal image, grid=major,
    every axis x label/.style={at={(ticklabel* cs:1)}, anchor=north east, yshift=-0.5mm, font=\sffamily\small, text=bkInk},
    every axis y label/.style={at={(ticklabel* cs:1)}, anchor=south, yshift=1mm, font=\sffamily\small, text=bkInk}},
  kixg sig/.style={kixgW, line width=1.4pt, no markers, samples=241, smooth},
  kixg part/.style={bkGreen!45!black, line width=0.8pt, dash pattern=on 3pt off 1.5pt, no markers, samples=241, smooth},
  kixg bar/.style={ycomb, draw=kixgF, line width=3pt, no markers},
  kixg dot/.style={only marks, mark=*, mark size=1.3pt, kixgW}}
\tikzset{
  kixg world/.style={draw, line width=1pt, rounded corners=1.5mm, fill=white, font=\sffamily\bfseries\small, text=bkInk,
    minimum width=30mm, minimum height=7mm, inner xsep=2mm},
  kixg bridge/.style={-{Stealth[length=2.4mm]}, line width=1pt, bkInk!80},
  kixg blab/.style={font=\sffamily\footnotesize, text=bkInk, inner sep=0.8mm}}
\newcommand\kixgworlds[5][42mm]{%
  \node[kixg world, draw=kixgW] (kixgL) at (23.5mm,#1) {#2};
  \node[kixg world, draw=kixgF] (kixgR) at (99.5mm,#1) {#3};
  \draw[kixg bridge] ([yshift=1.3mm, xshift=1.5mm]kixgL.east) -- node[kixg blab, above] {#4} ([yshift=1.3mm, xshift=-1.5mm]kixgR.west);
  \draw[kixg bridge] ([yshift=-1.3mm, xshift=-1.5mm]kixgR.west) -- node[kixg blab, below] {#5} ([yshift=-1.3mm, xshift=1.5mm]kixgL.east);}
\newcommand\kixgpole[3][kixgP]{\draw[#1, line width=1.3pt] (axis cs:#2,#3) ++(-1.3mm,-1.3mm) -- ++(2.6mm,2.6mm) (axis cs:#2,#3) ++(-1.3mm,1.3mm) -- ++(2.6mm,-2.6mm);}
\newcommand\kixgzero[3][kixgP]{\draw[#1, line width=1.1pt, fill=white] (axis cs:#2,#3) circle[radius=1.2mm];}
\makeatother
%% ===== この章の部品 =====
\makeatletter
% 小見出し（太字の1行）
\newcommand\kixghead[1]{\par\addvspace{6pt}\Needspace{6\baselineskip}\noindent{\sffamily\bfseries\color{\bkTT}#1}\par\nobreak\vspace{2pt}\noindent\ignorespaces}
% 囲みの中の小見出し（それだけで1行）
\newcommand\kixgbox[1]{\par\addvspace{3pt}\noindent\textbf{#1}\par\vspace{3pt}\noindent\ignorespaces}
% 式を、まん中に1行で置く
\newcommand\kixgline[1]{\par\addvspace{3pt}{\centering #1\par}\addvspace{3pt}\noindent\ignorespaces}
% 大事な式のわく  \kixgframe[はば]{中身}
\newcommand\kixgframe[2][100mm]{\par\addvspace{4pt}{\centering\tikz\node[draw=\bkTT, line width=0.9pt, fill=bkTheme!6, rounded corners=1.5mm, inner xsep=6mm, inner ysep=2.5mm, align=center]{#2};\par}\addvspace{4pt}\noindent\ignorespaces}
% 例題の中の式の変形の表（bkhenkei と同じ形で、行の間を少しつめたもの）
\newenvironment{kixghenkei}[1][0.48]{\par\addvspace{2pt}\noindent\renewcommand\arraystretch{1.22}%
  \begin{tabular}{@{\hspace{1\zw}}>{$\displaystyle}m{#1\linewidth}<{$}@{\hspace{2\zw}}>{\small\color{bkGray}\raggedright\arraybackslash}m{\dimexpr0.90\linewidth-#1\linewidth\relax}@{}}}%
  {\end{tabular}\par\addvspace{2pt}}
% 例題の途中式の2行め以降
\newcommand\kixga{\par\vspace{2pt}\noindent\rightskip=0pt plus 8em\relax\hangindent=2.6\zw\hangafter=0\relax}
\newcommand\kixgs{\rightskip=0pt plus 8em\relax}
\makeatother

%% 9-6 変換表と性質（第15回：組 V9-4。ページ制限なし。偶数ページ）
%% 「名前と操作」の型（tools/kata9-int.md）と、表計算・2つの世界の図の型（tools/kata9.md）で書く。
%% 場合分けの表9-6（A〜I。I たたみこみは この章で足した行）と流れ図、変換表の見開き（関数の表・性質の表）。
%% ほかの巻との約束：5-18:441「ラプラス変換の表（9-6）にならぶのも、この章の表と同じ関数」、6-13:664「たたみこみ（9-6・9-7）」、
%%   物理 S8（デルタ関数とインパルス応答）・S16（最終値の定理）、6-7 の δ（対数減衰率）とは別の文字の使い方、
%%   ステップ関数は H(t)（第15回の決定3）、たたみこみの積分の文字は t1（τ は時定数）。
%%   数はすべて Python で検算した（scratchpad/v9-4/check1.py〜check3.py）。

\bkchapter{9}{6}{変換表と性質}{\bklevel{5}}{変換表と性質}
\bkmichishirube{5-14,5-15,5-16,5-18,6-7,6-10,6-13,8-6,9-5}{9-7,9-8,9-9}

\begin{goal}
\begin{itemize}[label={\color{\bkTT}●}, leftmargin=1.6\zw, itemsep=0pt, topsep=1pt]
\item \bklevel{3}\enspace 表計算で $f(t)e^{-st}$ を台形でたし、変換表の値と同じになることを確かめられる。
\item \bklevel{4}\enspace $t^n$・$\sin$・$\cos$ と、$e^{-at}$ がかかった形の変換表の行を、部分積分・オイラーの公式・$s$ のずらしで作れる。
\item \bklevel{5}\enspace 場合分けの表9-6と流れ図を使って、おくれ・パルス・デルタ関数・たたみこみをふくむ関数を変換し、初期値・最終値の定理で確かめられる。
\end{itemize}
\end{goal}

\begin{hatena}
\noindent\begin{minipage}[c]{0.5\linewidth}
\raggedright
5-18 で、微分方程式の「答えの候補カード」を作りました。定数、$e^{at}$、$\sin\omega t$、$\cos\omega t$、$e^{at}\sin\omega t$、$t$ です。5-18 には「ラプラス変換の表（9-6）にならぶのも、この章の表と同じ関数です」と書きました。\par\vspace{4pt}
\textbf{候補カードの関数は、$s$ の世界では、どんな式に翻訳されるでしょうか。}
\end{minipage}\hfill
\begin{minipage}[c]{0.46\linewidth}
\centering
\begin{tikzpicture}[x=1mm, y=1mm, kixgc/.style={draw=bkBlue!80!black, line width=0.8pt, rounded corners=1mm, fill=white, minimum width=24mm, minimum height=6mm, font=\footnotesize, text=bkInk, inner sep=0.5mm}, kixgq/.style={draw=bkOrange!85!black, line width=0.8pt, rounded corners=1mm, fill=white, minimum width=16mm, minimum height=6mm, font=\sffamily\bfseries\footnotesize, text=bkInk}]
  \node[font=\sffamily\bfseries\footnotesize, text=\bkTT] at (12,52) {時間の世界};
  \node[font=\sffamily\bfseries\footnotesize, text=\bkTT] at (48,52) {$s$ の世界};
  \foreach \kixgy/\kixgf in {44/{定数 $C$},36/{$e^{at}$},28/{$\sin\omega t$},20/{$\cos\omega t$},12/{$e^{at}\sin\omega t$},4/{$t$}} {
    \node[kixgc] at (12,\kixgy) {\kixgf};
    \draw[bkInk!70, line width=0.7pt, -{Stealth[length=1.6mm]}] (25.5,\kixgy) -- (38.5,\kixgy);
    \node[kixgq] at (48,\kixgy) {?};}
\end{tikzpicture}
\end{minipage}
\end{hatena}

\noindent 9-5 では、翻訳の辞書（変換表）の行を、定数・$t$・$e^{at}$ と、微分の行まで作りました。この章では、機械や回路の動きに出てくる関数の行を全部そろえて、変換表を完成させます。表の行は、どれも1回だけ自分で作り、あとは引くだけにします。行を作る道具は、第5巻の部分積分、第6巻のオイラーの公式、そしてこの章で作る「ずらし」の3つだけです。

表を引くだけでは足りない形もあります。スイッチをおくれて入れる、短い間だけ電圧をかける、ハンマーでたたく、ずっと入れ続けたものの効き目をたし合わせる、などです。これらは、表の行どうしを組み合わせる「性質」で計算します。どの場合にどうするかを、場合分けの表と流れ図にまとめます。行の数が多い章ですが、1回で全部覚えなくても大丈夫です。表は、引くためにあります。

\begin{omoidasou}
\bkrecall{9-5}{$\mathcal L[f]=\displaystyle\int_0^\infty f(t)e^{-st}\,\dd t$。$\mathcal L[1]=\tfrac{1}{s}$、$\mathcal L[t]=\tfrac{1}{s^2}$、$\mathcal L[e^{at}]=\tfrac{1}{s-a}$、$\mathcal L[f']=sF-f(0)$。}
\bkrecall{5-16}{部分積分 $\displaystyle\int f\,g'\,\dd t=f\,g-\int f'\,g\,\dd t$。}
\bkrecall{6-10}{オイラーの公式 $e^{j\theta}=\cos\theta+j\sin\theta$。複素数の実部と虚部を分けて読む。}
\bkrecall{5-15}{$(\sin\omega t)'=\omega\cos\omega t$、$(\cos\omega t)'=-\omega\sin\omega t$。2回微分すると $-\omega^2$ 倍。}
\bkrecall{5-16}{置換積分：$u=t-t_0$ と置くと $\dd u=\dd t$。積分の範囲も $u$ の値に直す。}
\bkrecall{8-6}{重積分は、内側と外側の積分の順番を入れかえてよい。三角形の範囲は、たて線でも横線でも切れる（場合B）。}
\end{omoidasou}

%% ---------------- やってみよう ----------------
\begin{yattemiyou}[表計算で、4つの関数を翻訳しよう]
$s=2$ として、$1$、$t$、$e^{-t}$、$\sin3t$ の4つの関数に $e^{-2t}$ をかけ、$t=0$ から10まで台形でたします（9-5 の表計算の続き）。
\begin{enumerate}[label={\bkcircled[bkGreen]{\arabic*}}, leftmargin=2.4\zw, itemsep=1pt, topsep=2pt]
\item A列に時刻 $t$ を、0から10まで0.1ずつ入れる（A2〜A102）。
\item B2 に \texttt{=EXP(-2*A2)} を入れる（かけるもの $e^{-2t}$）。
\item C2〜F2 に、$1\times e^{-2t}$、$t\,e^{-2t}$、$e^{-t}e^{-2t}$、$\sin3t\,e^{-2t}$ を入れる（\texttt{=B2}、\texttt{=A2*B2}、\texttt{=EXP(-A2)*B2}、\texttt{=SIN(3*A2)*B2}）。
\item 2行めを102行めまでコピーする。
\item 103行めに台形の和 \texttt{=0.1*(SUM(C2:C102)-(C2+C102)/2)} を入れ、右へコピーする。
\end{enumerate}
\end{yattemiyou}

\noindent ⑤の式は、9-5 の表の「台形」の列をたしたものと同じです。台形をすべてたすと、両はしの値だけが半分の重みになるので、「全部の合計から、両はしの半分をひいて、$\Delta t=0.1$ をかける」になります。
\begin{bkcenter}
\bksheet{6}{16}{%
1 & \SetCell{c} $t$ & \SetCell{c} $e^{-2t}$ & \SetCell{bg=kixgN} $1$ & \SetCell{bg=kixgN} $t$ & \SetCell{bg=kixgN} $e^{-t}$ & \SetCell{bg=kixgN} $\sin3t$ \\
2 & 0 & 1.0000 & 1.0000 & 0.0000 & 1.0000 & 0.0000 \\
3 & 0.1 & 0.8187 & 0.8187 & 0.0819 & 0.7408 & 0.2420 \\
4 & 0.2 & 0.6703 & 0.6703 & 0.1341 & 0.5488 & 0.3785 \\
5 & 0.3 & 0.5488 & 0.5488 & 0.1646 & 0.4066 & 0.4299 \\
$\vdots$ & $\vdots$ & $\vdots$ & $\vdots$ & $\vdots$ & $\vdots$ & $\vdots$ \\
102 & 10 & 0.0000 & 0.0000 & 0.0000 & 0.0000 & 0.0000 \\
103 & \SetCell{c} 台形 & & \SetCell{bg=kixgS} 0.50167 & \SetCell{bg=kixgS} 0.24917 & \SetCell{bg=kixgS} 0.33583 & \SetCell{bg=kixgS} 0.22827 \\}
\bktabcaption{4つの関数に $e^{-2t}$ をかけた表（小数第4位まで）。1行めは、かける前の関数。103行めが、$s=2$ のときのラプラス変換の近い値。}
\label{bk:9-6:tab-sheet}
\end{bkcenter}

%% ---------------- 数学の道具 ----------------
\Needspace{14\baselineskip}\bksection{doguu}{数学の道具}

\bkstage{me}{目でわかる}{表計算の値と、表の式をくらべる}{\bklevel{3}}

\bksubsection{1}{面積が、変換の値}

\label{bk:9-6:menseki}表計算の103行めは、$f(t)e^{-2t}$ のグラフの下の面積です（下の面積は正、上の面積は負：5-10）。グラフで見ると、$e^{-2t}$ をかけたおかげで、どの関数も $t=3$ くらいでほとんど0になり、面積が決まった数になっています。

\begin{bkfigure}
\begin{tikzpicture}
\begin{axis}[bkplot, width=112mm, height=52mm, xmin=0, xmax=4, ymin=-0.1, ymax=1.15, xtick={1,2,3,4}, ytick={0.5,1}, xlabel={$t$}, ylabel={$f(t)\,e^{-2t}$}]
  \addplot[bkcurve, domain=0:3.9, samples=80] {exp(-2*x)};
  \addplot[bkcurve2, dashed, domain=0:3.9, samples=80] {x*exp(-2*x)};
  \addplot[bkcurve3, dash pattern=on 1pt off 1.5pt, domain=0:3.9, samples=80] {exp(-3*x)};
  \addplot[bkInk, line width=1.2pt, dash pattern=on 5pt off 1.5pt on 1pt off 1.5pt, domain=0:3.9, samples=160] {sin(deg(3*x))*exp(-2*x)};
  \node[font=\sffamily\footnotesize, anchor=west] at (axis cs:1.5,1.0) {青の線：$1\times e^{-2t}$（面積 0.5）};
  \node[font=\sffamily\footnotesize, anchor=west] at (axis cs:1.5,0.86) {オレンジの破線：$t\,e^{-2t}$（面積 0.25）};
  \node[font=\sffamily\footnotesize, anchor=west] at (axis cs:1.5,0.72) {緑の点線：$e^{-t}e^{-2t}$（面積 0.333）};
  \node[font=\sffamily\footnotesize, anchor=west] at (axis cs:1.5,0.58) {黒の一点鎖線：$\sin3t\,e^{-2t}$（面積 0.231）};
\end{axis}
\end{tikzpicture}
\bkfigcaption{4つの関数に $e^{-2t}$ をかけたグラフ。下の面積が、$s=2$ のときのラプラス変換。$\sin3t\,e^{-2t}$ は、$t=1.05$ 〜 2.09 などで負になり、その分の面積はひかれる。}
\label{bk:9-6:fig-menseki}
\end{bkfigure}

\noindent この章で作る表の式に $s=2$ を入れた値と、表計算の値をくらべます。

\begin{bkcenter}
\begin{bktable}{colspec={Q[c,m,22mm] Q[c,m,26mm] Q[c,m,18mm] Q[c,m,22mm] Q[c,m,22mm]}, row{2-Z}={font=\small}, rowsep=2pt}
$f(t)$ & 表の式 $F(s)$ & $s=2$ & 台形 $\Delta t=0.1$ & 台形 $\Delta t=0.01$ \\
$1$ & $\frac{1}{s}$ & 0.5 & 0.50167 & 0.50002 \\
$t$ & $\frac{1}{s^2}$ & 0.25 & 0.24917 & 0.24999 \\
$e^{-t}$ & $\frac{1}{s+1}$ & 0.33333 & 0.33583 & 0.33336 \\
$\sin3t$ & $\frac{3}{s^2+9}$ & 0.23077 & 0.22827 & 0.23074 \\
\end{bktable}
\bktabcaption{表の式（この章で作る）と、表計算の値。$\Delta t$ を小さくすると、表の式の値に近づく。}
\label{bk:9-6:tab-kurabe}
\end{bkcenter}
$\Delta t$ を $\frac{1}{10}$ にすると、ずれはおよそ $\frac{1}{100}$ になります（台形の和のずれは、$\Delta t$ の2乗に比例して小さくなる：5-10 の考えの続き）。表の式は、$\Delta t$ を限りなく小さくした先の値を、1回の計算で出したものです。

\bkstage{te}{手でわかる}{表の行を、自分で作る}{\bklevel{4}}

\bksubsection{2}{$t^n$ の行：部分積分をくり返す}

\label{bk:9-6:tn}\kixgdanhead{1}{記号を読む：$\mathcal L[t^n]$}
$\mathcal L[t^n]$ は、「$t^n$ に $e^{-st}$ をかけて、0から $\infty$ までたす」の名前です。名前のままでは数は出ません。1回だけ、ふつうの定積分（部分積分）に直して、行を作ります。
\[ \kixgNt{\mathcal L[t^n]}=\kixgSt{1}{\int_0^{\infty}t^n\,e^{-st}\,\dd t} \]
\kixgdanhead{2}{部分積分で、$n$ を1つずつ下げる}
$t^n$ を $f$（微分するほう）、$e^{-st}$ を $g'$（積分するほう、$g=-\frac{1}{s}e^{-st}$）にします。
\begin{bkhenkei}[0.6]
\int_0^{\infty}t^ne^{-st}\,\dd t & 求めたいもの \\
=\Bigl[-\frac{t^n}{s}e^{-st}\Bigr]_0^{\infty}+\frac{n}{s}\int_0^{\infty}t^{n-1}e^{-st}\,\dd t & 部分積分。$(t^n)'=nt^{n-1}$ \\
=0+\frac{n}{s}\,\mathcal L[t^{n-1}] & はしの項は0（$t=0$ で0、$\infty$ 側は 5-18） \\
\end{bkhenkei}
\kixgframe[100mm]{\Large $\mathcal L[t^n]=\dfrac{n}{s}\,\mathcal L[t^{n-1}]$}
「$t$ の次数が1つ下がり、$\frac{n}{s}$ がかかる」というきまりです。$\mathcal L[1]=\frac{1}{s}$ から始めて、くり返します。
\begin{bkhenkei}[0.5]
\mathcal L[t]=\frac{1}{s}\times\frac{1}{s}=\frac{1}{s^2} & $n=1$（9-5 と同じ） \\
\mathcal L[t^2]=\frac{2}{s}\times\frac{1}{s^2}=\frac{2}{s^3} & $n=2$ \\
\mathcal L[t^3]=\frac{3}{s}\times\frac{2}{s^3}=\frac{6}{s^4}=\frac{3!}{s^4} & $n=3$。$3\times2\times1=3!$ \\
\end{bkhenkei}
分子は $1$、$2\times1$、$3\times2\times1$ と、かけ算が1つずつふえていきます。だから $\mathcal L[t^n]=\dfrac{n!}{s^{n+1}}$ です（$n!$ は5-17 の階乗）。表計算の $t$ の列（$s=2$）は $\frac{1}{4}=0.25$ に近い値でした。

\bksubsection{3}{$\sin$ と $\cos$ の行：2通りで作る}

\label{bk:9-6:sincos}\kixghead{1つめの作り方：オイラーの公式}
9-5 の $e^{at}$ の行 $\frac{1}{s-a}$ は、$a$ が複素数でも同じ計算で成り立ちます（$s$ の実部が $a$ の実部より大きいとき）。$a=j\omega$ とすると、$e^{j\omega t}=\cos\omega t+j\sin\omega t$（6-10）なので、
\begin{bkhenkei}[0.52]
\mathcal L[e^{j\omega t}]=\frac{1}{s-j\omega} & $e^{at}$ の行で $a=j\omega$ \\
=\frac{s+j\omega}{(s-j\omega)(s+j\omega)}=\frac{s+j\omega}{s^2+\omega^2} & 分母・分子に $s+j\omega$ をかける（$j^2=-1$） \\
=\frac{s}{s^2+\omega^2}+j\,\frac{\omega}{s^2+\omega^2} & 実部と虚部に分ける \\
\mathcal L[\cos\omega t]+j\,\mathcal L[\sin\omega t] & 左辺は、線形で2つに分ける \\
\end{bkhenkei}
$s$ が実数なら、実部どうし・虚部どうしが等しいので、
\kixgframe[110mm]{\Large $\mathcal L[\cos\omega t]=\dfrac{s}{s^2+\omega^2}$\qquad $\mathcal L[\sin\omega t]=\dfrac{\omega}{s^2+\omega^2}$}

\kixghead{2つめの作り方：微分の行で確かめる}
$f=\sin\omega t$ は、2回微分すると $-\omega^2$ 倍になり（5-15）、$f(0)=0$、$f'(0)=\omega$ です。9-5 の2回微分の行を使います。
\begin{bkhenkei}[0.5]
\mathcal L[f'']=-\omega^2F & $f''=-\omega^2f$ の両辺を変換 \\
s^2F-s\cdot0-\omega=-\omega^2F & 左辺に2回微分の行 \\
(s^2+\omega^2)F=\omega,\quad F=\frac{\omega}{s^2+\omega^2} & $F$ についてまとめる \\
\end{bkhenkei}
同じ式が出ました。$\cos\omega t$ なら $f(0)=1$、$f'(0)=0$ で、$s^2F-s=-\omega^2F$ から $F=\frac{s}{s^2+\omega^2}$ です。$s=2$、$\omega=3$ では $\mathcal L[\sin3t]=\frac{3}{4+9}=0.23077$ で、表計算の値 0.22827（$\Delta t=0.01$ で 0.23074）に合っています。

\bksubsection{4}{$s$ のずらし：$e^{-at}$ がかかった形}

\label{bk:9-6:szurashi}5-18 のカードの「ゆれながら止まる」形 $e^{-at}\sin\omega t$ は、$\sin\omega t$ に $e^{-at}$ がかかった形です。定義の積分の中で、$e^{-at}$ と $e^{-st}$ をまとめます。
\begin{bkhenkei}[0.56]
\mathcal L[e^{-at}f(t)]=\int_0^{\infty}f(t)\,e^{-at}e^{-st}\,\dd t & 名前を、定義の積分に \\
=\int_0^{\infty}f(t)\,e^{-(s+a)t}\,\dd t & 肩のたし算（5-14） \\
=F(s+a) & $F(s)$ の式の $s$ を、$s+a$ に取りかえたもの \\
\end{bkhenkei}
\kixgframe[100mm]{\Large $\mathcal L[e^{-at}f(t)]=F(s+a)$}
$e^{-at}$ をかけることは、$s$ の世界では「$s$ を $a$ だけ大きくする（ずらす）」ことです。これを\textbf{$s$ のずらし}といいます。$F(s)$ を求めてから、$s$ を $s+a$ に書きかえます。
\begin{bkhenkei}[0.5]
\mathcal L[e^{-at}\sin\omega t]=\frac{\omega}{(s+a)^2+\omega^2} & $\frac{\omega}{s^2+\omega^2}$ の $s$ を $s+a$ に \\
\mathcal L[e^{-at}\cos\omega t]=\frac{s+a}{(s+a)^2+\omega^2} & $\frac{s}{s^2+\omega^2}$ の $s$ を $s+a$ に（分子の $s$ も） \\
\mathcal L[t\,e^{-at}]=\frac{1}{(s+a)^2} & $\frac{1}{s^2}$ の $s$ を $s+a$ に \\
\end{bkhenkei}

\begin{bkfigure}
\begin{tikzpicture}
\kixgworlds[63mm]{時間の世界}{$s$ の世界}{ラプラス変換}{逆変換}
\begin{axis}[kixg wave, at={(0mm,0mm)}, height=47mm, xmin=0, xmax=5, ymin=-1.15, ymax=1.15, xtick={0,1,...,5}, ytick={-1,1}, xlabel={時間 $t$}, ylabel={$x$}]
  \addplot[kixgF, line width=1.2pt, dash pattern=on 4pt off 2pt, no markers, samples=241, smooth, domain=0:5] {sin(deg(2*x))};
  \addplot[kixg sig, domain=0:5] {exp(-x)*sin(deg(2*x))};
\end{axis}
\begin{axis}[kixg splane, at={(76mm,0mm)}, xmin=-3, xmax=2, ymin=-2.5, ymax=2.5, xtick={-2,-1,1}, ytick={-2,-1,1,2}, yticklabels={,$-1$,$1$,}, xticklabel style={yshift=-1.4mm}, xlabel={実部}, ylabel={虚部}]
  \node[font=\footnotesize, anchor=west] at (axis cs:0.22,2) {$2$};
  \node[font=\footnotesize, anchor=west] at (axis cs:0.22,-2) {$-2$};
  \kixgpole[kixgF]{0}{2}\kixgpole[kixgF]{0}{-2}
  \kixgpole[kixgW]{-1}{2}\kixgpole[kixgW]{-1}{-2}
  \draw[bkInk!75, line width=0.9pt, -{Stealth[length=2mm]}] (axis cs:-0.2,1.45) -- (axis cs:-0.8,1.45);
  \node[font=\sffamily\footnotesize, anchor=east, fill=white, inner sep=1pt] at (axis cs:-1.15,1.45) {左へ1};
\end{axis}
\end{tikzpicture}
\bkfigcaption{左：$\sin2t$（オレンジの破線）と $e^{-t}\sin2t$（青い線）。右：それぞれの変換の分母が0になる所（オレンジの×が $\sin2t$、青の×が $e^{-t}\sin2t$）。$e^{-t}$ をかけると、×が左へ1だけずれる。}
\label{bk:9-6:fig-zurashi}
\end{bkfigure}

\noindent $\sin2t$ の変換 $\frac{2}{s^2+4}$ は、$s=\pm2j$ で分母が0です。$e^{-t}\sin2t$ の変換 $\frac{2}{(s+1)^2+4}$ は、$s+1=\pm2j$、つまり $s=-1\pm2j$ で分母が0です。分母が0になる点が、そろって左へ1だけずれました。「ずっとゆれる」動きが「ゆれながら止まる」動きに変わると、点は左へ動くのです（9-9 の極の読み方の入口）。車B の $s^2+2s+5$ は $(s+1)^2+4$ と書けるので、この形です（9-7）。

\bkkokomade[表計算の値と表の式をくらべ、$t^n$・$\sin$・$\cos$ の行を部分積分とオイラーの公式で作り、$e^{-at}$ がかかった形を $s$ のずらしで表を引けることがわかれば十分です。この先は、おくれ・パルス・デルタ関数・たたみこみ・最終値の定理と、それらをまとめた場合分けの表です。]

\Needspace{9\baselineskip}\bkstage{shiki}{式でわかる}{性質の行と、場合分けの表}{\bklevel{5}}

\bksubsection{5}{時間のずらし：おくれて始まる（場合 D・F）}

\label{bk:9-6:jikan}\kixghead{スイッチを $t_0$ 秒後に入れる}
9-5 のステップ関数 $H(t)$ は、時刻0にスイッチを入れる形でした。$t_0$ 秒後に入れるなら、グラフを右へ $t_0$ ずらした $H(t-t_0)$ です（3-9 の放物線を動かす考え：$x$ を $x-p$ にすると右へ $p$）。一般に、$f(t)$ のグラフを右へ $t_0$ ずらし、$t<t_0$ では0にした関数は $f(t-t_0)H(t-t_0)$ と書けます。

\begin{bkfigure}
\begin{tikzpicture}
\begin{axis}[bkplot, width=112mm, height=40mm, xmin=0, xmax=5, ymin=0, ymax=1.25, xtick={1,2,3,4,5}, ytick={0.5,1}, xlabel={$t$}]
  \addplot[bkthin, dashed, domain=0:5, samples=60] {exp(-x)};
  \addplot[bkcurve, domain=1:5, samples=60] {exp(-(x-1))};
  \addplot[bkcurve, domain=0:1] {0};
  \draw[bkBlue, line width=0.8pt, dash pattern=on 1.5pt off 1pt] (axis cs:1,0) -- (axis cs:1,1);
  \fill[bkBlue] (axis cs:1,1) circle (1.6pt);
  \draw[bkInk!75, line width=0.9pt, -{Stealth[length=2mm]}] (axis cs:0.3,1.08) -- (axis cs:0.92,1.08);
  \node[font=\sffamily\footnotesize, anchor=south west] at (axis cs:0.25,1.11) {1秒おくれる};
  \node[font=\sffamily\footnotesize, anchor=west, fill=white, inner sep=1pt] at (axis cs:2.2,0.8) {灰色の破線：$e^{-t}$};
  \node[font=\sffamily\footnotesize, anchor=west, fill=white, inner sep=1pt] at (axis cs:2.2,0.52) {青い線：$e^{-(t-1)}H(t-1)$};
\end{axis}
\end{tikzpicture}
\bkfigcaption{$e^{-t}$（灰色の破線）を右へ1ずらし、$t<1$ を0にした関数（青い線）。}
\label{bk:9-6:fig-okure}
\end{bkfigure}

\kixghead{変換：$e^{-st_0}$ をかける}
$t<t_0$ では関数が0なので、積分は $t_0$ からです。$u=t-t_0$ と置いて、置換積分（5-16）をします。
\begin{bkhenkei}[0.56]
\mathcal L[f(t-t_0)H(t-t_0)] & 求めたいもの \\
=\int_{t_0}^{\infty}f(t-t_0)\,e^{-st}\,\dd t & $t<t_0$ は0なので、$t_0$ から \\
=\int_0^{\infty}f(u)\,e^{-s(u+t_0)}\,\dd u & $u=t-t_0$。$t=t_0$ で $u=0$ \\
=e^{-st_0}\int_0^{\infty}f(u)\,e^{-su}\,\dd u=e^{-st_0}F(s) & $e^{-st_0}$ は $u$ に関係しないので外へ \\
\end{bkhenkei}
\kixgframe[110mm]{\Large $\mathcal L[f(t-t_0)\,H(t-t_0)]=e^{-st_0}\,F(s)$}
時間の世界で「$t_0$ おくれる」ことは、$s$ の世界では「$e^{-st_0}$ をかける」ことです。これを\textbf{時間のずらし}といいます。たとえば $\mathcal L[H(t-2)]=\frac{e^{-2s}}{s}$、$\mathcal L[e^{-(t-1)}H(t-1)]=\frac{e^{-s}}{s+1}$ です。$s$ のずらし（④）と向きがちがうので注意します：$e^{-at}$ を「かける」と $s$ がずれ、時間を「ずらす」と $e^{-st_0}$ がかかります。

\kixghead{パルス：ステップ関数の和に直す（場合 F）}
はじめの1秒間だけ1\,V をかけ、そのあと0\,V にする電圧（幅1秒・高さ1の\textbf{パルス}）は、区間ごとに式が変わります。「時刻0に1を足し、時刻1に1をひく」と考えると、
\kixgline{パルス $=H(t)-H(t-1)$}
と、ステップ関数の和に直せます。あとは項ごとに表と時間のずらしで、
\kixgline{$\mathcal L[H(t)-H(t-1)]=\dfrac{1}{s}-\dfrac{e^{-s}}{s}=\dfrac{1-e^{-s}}{s}$}
です。$s=2$ では $\frac{1-e^{-2}}{2}=0.43233$ です。

\begin{bkfigure}
\begin{tikzpicture}[x=1mm, y=1mm]
  \foreach \kixgx/\kixgt in {0/{$H(t)$},46/{$-H(t-1)$},92/{パルス}} {
    \draw[bkInk, line width=0.6pt, -{Stealth[length=1.6mm]}] (\kixgx,0) -- (\kixgx+38,0) node[right, font=\footnotesize] {$t$};
    \draw[bkInk, line width=0.6pt, -{Stealth[length=1.6mm]}] (\kixgx,-12) -- (\kixgx,14);
    \node[font=\footnotesize, anchor=south] at (\kixgx+18,15) {\kixgt};
    \node[font=\footnotesize, anchor=north east] at (\kixgx+11.4,-1) {1};
    \draw[bkInk, line width=0.5pt] (\kixgx+12,0) -- (\kixgx+12,-1);}
  \draw[bkBlue, line width=1.6pt] (0,10) -- (34,10);
  \draw[bkOrange!85!black, line width=1.6pt] (46,0) -- (58,0) (58,-10) -- (80,-10);
  \draw[bkOrange!85!black, line width=0.8pt, dash pattern=on 1.5pt off 1pt] (58,0) -- (58,-10);
  \draw[bkGreen!50!black, line width=1.6pt] (92,10) -- (104,10) (104,0) -- (126,0);
  \draw[bkGreen!50!black, line width=0.8pt, dash pattern=on 1.5pt off 1pt] (104,10) -- (104,0);
  \node[font=\footnotesize, anchor=east] at (-1,10) {1};
  \node[font=\footnotesize, anchor=east] at (45,-10) {$-1$};
  \node[font=\sffamily\bfseries\large] at (42,9) {＋};
  \node[font=\sffamily\bfseries\large] at (88,9) {＝};
\end{tikzpicture}
\bkfigcaption{パルスは、時刻0に始まるステップ（左）と、時刻1に始まる $-1$ のステップ（まん中）の和（右）。}
\label{bk:9-6:fig-pulse}
\end{bkfigure}

\kixghead{むだ時間}
お風呂のシャワーの温度を変えても、新しいお湯が配管を通って出てくるまで、少し時間がかかります。このように「入れたものが、そのまま $L$ 秒おくれて出てくる」ことを\textbf{むだ時間}といいます。時間のずらしから、むだ時間は $s$ の世界では $e^{-sL}$ をかけることです。ベルトコンベアや、長い配管の工場の制御で大切になります（物理 波-4）。

\bksubsection{6}{積分の行と、$t$ をかける行（場合 E・G）}

\label{bk:9-6:tkakeru}\kixghead{積分の行：$s$ でわる}
$g(t)=\displaystyle\int_0^tf(t_1)\,\dd t_1$（0から $t$ までの面積）とすると、$g'=f$、$g(0)=0$ です（5-11）。微分の行を $g$ に使います。
\begin{bkhenkei}[0.5]
\mathcal L[g']=sG(s)-g(0) & 微分の行 \\
F(s)=sG(s)-0 & $g'=f$、$g(0)=0$ \\
G(s)=\frac{F(s)}{s} & $s$ でわる \\
\end{bkhenkei}
微分は「$s$ をかける」、積分は「$s$ でわる」です。$f=1$ なら $\int_0^t1\,\dd t_1=t$ で、$\frac{1}{s}\div s=\frac{1}{s^2}$ は $t$ の行と合っています。

\kixghead{$t$ をかける行：$s$ で微分する}
$F(s)=\int_0^{\infty}f(t)e^{-st}\,\dd t$ を $s$ で微分します。積分の中の $e^{-st}$ を $s$ で微分すると、$-te^{-st}$ です（$t$ は止めておく：8-2 の偏微分と同じ考え）。
\begin{bkhenkei}[0.56]
\frac{\dd F}{\dd s}=\int_0^{\infty}f(t)\,(-t)\,e^{-st}\,\dd t & 積分の中で $s$ で微分 \\
-\frac{\dd F}{\dd s}=\int_0^{\infty}t\,f(t)\,e^{-st}\,\dd t=\mathcal L[t\,f(t)] & 両辺に $-1$ をかける \\
\end{bkhenkei}
\kixgframe[100mm]{\Large $\mathcal L[t\,f(t)]=-\dfrac{\dd F}{\dd s}$}
\textbf{2通りで確かめる}：$te^{-t}$ は、$s$ のずらし（$\frac{1}{s^2}$ の $s$ を $s+1$ に）で $\frac{1}{(s+1)^2}$。$t$ をかける行でも、$F=\frac{1}{s+1}$ を微分して $-\frac{\dd F}{\dd s}=\frac{1}{(s+1)^2}$。同じです。この形は、9-7 の「重解」で使います。表の $t\sin\omega t$・$t\cos\omega t$ の行も、$\sin\omega t$・$\cos\omega t$ の行をこの方法で微分して作ったものです（$\cos$ の行は、商の微分（5-6）を使う）。

\bksubsection{7}{デルタ関数：たたいた力を表す（場合 A）}

\label{bk:9-6:delta}ハンマーで機械をコツンとたたくと、力はとても短い時間だけはたらきます。大事なのは力の大きさそのものより、「力×時間」（力積。物理 力-3）です。そこで、面積が1の細いパルスを考えます。

\begin{bkfigure}
\begin{tikzpicture}[x=1mm, y=1mm]
  \draw[bkInk, line width=0.6pt, -{Stealth[length=1.8mm]}] (-2,0) -- (62,0) node[right, font=\footnotesize] {$t$};
  \draw[bkInk, line width=0.6pt, -{Stealth[length=1.8mm]}] (0,-2) -- (0,46);
  \draw[bkBlue, line width=1.2pt] (0,0) -- (0,10) -- (40,10) -- (40,0);
  \draw[bkOrange!85!black, line width=1.2pt, dash pattern=on 4pt off 1.5pt] (0,20) -- (20,20) -- (20,0);
  \draw[bkGreen!50!black, line width=1.2pt, dash pattern=on 1pt off 1.2pt] (0,40) -- (10,40) -- (10,0);
  \foreach \kixgy/\kixgl in {10/1,20/2,40/4} {\node[font=\footnotesize, anchor=east] at (-1,\kixgy) {\kixgl};}
  \foreach \kixgx/\kixgl in {10/{$\frac14$},20/{$\frac12$},40/1} {\node[font=\footnotesize, anchor=north] at (\kixgx,-1) {\kixgl};}
  \node[font=\footnotesize, anchor=west, align=left] at (66,38) {緑の点線：幅 $\frac14$・高さ4};
  \node[font=\footnotesize, anchor=west, align=left] at (66,30) {オレンジの破線：幅 $\frac12$・高さ2};
  \node[font=\footnotesize, anchor=west, align=left] at (66,22) {青い線：幅1・高さ1};
  \node[font=\footnotesize, anchor=west, align=left] at (66,12) {どれも面積は1};
\end{tikzpicture}
\bkfigcaption{面積が1のまま、幅を細くしたパルス。幅 $\varepsilon$、高さ $\frac{1}{\varepsilon}$。$\varepsilon$ を0に近づけた先がデルタ関数 $\delta(t)$。}
\label{bk:9-6:fig-delta}
\end{bkfigure}

幅 $\varepsilon$（イプシロン）・高さ $\frac{1}{\varepsilon}$ のパルスは、⑤のパルスの $\frac{1}{\varepsilon}$ 倍で、幅が $\varepsilon$ なので、変換は $\dfrac{1-e^{-\varepsilon s}}{\varepsilon s}$ です。$s=2$ で、幅を細くしていくと、次のようになります。
\begin{bkcenter}
\begin{bktable}{colspec={Q[c,m,30mm] Q[c,m,18mm] Q[c,m,18mm] Q[c,m,18mm] Q[c,m,18mm]}, row{2-Z}={font=\small}, rowsep=2pt}
幅 $\varepsilon$ & 1 & 0.1 & 0.01 & 0.001 \\
変換の値（$s=2$） & 0.432 & 0.906 & 0.990 & 0.999 \\
\end{bktable}
\bktabcaption{細いパルスの変換は、1に近づく（$e^{-\varepsilon s}\approx1-\varepsilon s$ だから：5-17 のテイラー展開の1次まで）。}
\end{bkcenter}
幅を限りなく細くした先を\textbf{デルタ関数}（でるたかんすう。ディラックのデルタ関数）といい、$\delta(t)$ と書きます。\bkidx{でるたかんすう}{デルタ関数}変換は、どの $s$ でも1です。
\kixgframe[90mm]{\Large $\mathcal L[\delta(t)]=1$}
「時刻0に、力積1でたたく」ことを表します。ここの $\delta$ は、6-7 の対数減衰率 $\delta$ とは別の文字の使い方です。力積 $P$ でたたくなら $P\,\delta(t)$、$t_0$ 秒後にたたくなら $\delta(t-t_0)$ で、時間のずらしから $\mathcal L[\delta(t-t_0)]=e^{-st_0}$ です。9-4 のせいかくに言うとで「くり返す波のスペクトルは、とびとびの線になる」と書いた、その「線」も、このデルタ関数です（周波数の世界で、幅は限りなく細く、面積だけが決まっている）。

\kixghead{インパルス応答}
機械や回路を $\delta(t)$ でたたいたときの動きを\textbf{インパルス応答}といいます。\bkidx{いんぱるすおうとう}{インパルス応答}たとえば $x'+x=\delta(t)$（はじめ $x=0$）なら、$sX+X=1$ から $X=\frac{1}{s+1}$ で、$x=e^{-t}$ です。たたいた瞬間に $x$ が1にとび上がり、そのあと $e^{-t}$ で小さくなっていきます。インパルス応答がわかると、どんな入力にもこたえられることを、次の⑧で見ます。

\Needspace{10\baselineskip}
\begin{seikaku}
\kixgbox{デルタ関数は、ふつうの関数ではない}
$\delta(t)$ は、$t\neq0$ ではどこでも0で、$t=0$ では「限りなく大きい」ものです。値で決まる関数ではなく、「面積（積分）が1」ということで決まるものです。だから、$\delta(t)$ の値を表に書いたり、グラフを高さでかいたりはできません（図では、高さ1の矢印でかく本が多い）。また、$\delta(t)$ は時刻0にあるので、ラプラス変換の積分は、0のほんの少し前（$0-$ と書く）から始める約束です。こうしないと、たたいた分が積分に入りません。微分の行の $f(0)$ も、そのときは $f(0-)$（たたく直前の値）です。数学として正しくあつかう方法（超関数）は、1945年ごろにシュワルツが作りました。
\end{seikaku}

\begin{motto}[5]
\kixgbox{デルタ関数のスペクトルは、どの周波数も同じ高さ}
9-4 で、手をたたく音のような短い音は、スペクトルが広い帯になることを見ました。9-4 の方形パルスの変換を使うと、まん中が $t=0$ で幅 $\varepsilon$・高さ $\frac{1}{\varepsilon}$ のパルスのフーリエ変換は $F(\omega)=\sin(\omega\varepsilon/2)\div(\omega\varepsilon/2)$ です。$\omega=0$ ではいつも1（面積）で、はじめて0になる $\omega$ は $\frac{2\pi}{\varepsilon}$ です。幅を細くするほど、1に近い所が右へ広がり（図\ref{bk:9-6:fig-deltaspec}）、デルタ関数では、どの $\omega$ でも1になります。ハンマーでたたいて機械を調べる（物理・工学ではこう使う）のは、1回たたくだけで、すべての周波数で機械をゆすったのと同じになるからです。
\end{motto}
\begin{bkfigure}
\begin{tikzpicture}
\kixgworlds{時間の世界}{周波数の世界}{変換する}{もどす}
\begin{axis}[kixg wave, at={(0mm,0mm)}, xmin=-1, xmax=1, ymin=0, ymax=4.6, xtick={-0.5,0,0.5}, ytick={1,4}, extra y tick labels={}, xlabel={時間 $t$}, ylabel={高さ}]
  \addplot[kixgW, line width=1.4pt, no markers] coordinates {(-1,0) (-0.5,0) (-0.5,1) (0.5,1) (0.5,0) (1,0)};
  \addplot[kixgF, line width=1.4pt, dash pattern=on 4pt off 2pt, no markers] coordinates {(-1,0) (-0.125,0) (-0.125,4) (0.125,4) (0.125,0) (1,0)};
\end{axis}
\begin{axis}[kixg wave, at={(76mm,0mm)}, xmin=0, xmax=30, ymin=-0.3, ymax=1.15, xtick={0,10,20,30}, ytick={1}, xlabel={角周波数 $\omega$（rad/s）}, ylabel={$F(\omega)$}]
  \addplot[kixg sig, domain=0.01:30] {sin(deg(x/2))/(x/2)};
  \addplot[kixgF, line width=1.4pt, dash pattern=on 4pt off 2pt, no markers, samples=121, smooth, domain=0.01:30] {sin(deg(x/8))/(x/8)};
\end{axis}
\end{tikzpicture}
\bkfigcaption{面積1のパルス（左）と、そのフーリエ変換（右）。青い線は幅1・高さ1（最初の0は $\omega=2\pi\approx6.3$）、オレンジの破線は幅 $\frac14$・高さ4（最初の0は $\omega=8\pi\approx25$）。細いほど、1に近い所が広い。}
\label{bk:9-6:fig-deltaspec}
\end{bkfigure}

\bksubsection{8}{たたみこみ：たたいた答えを、たし合わせる（場合 I）}

\label{bk:9-6:tatami}6-13 では、「$\cos$ でない一般の形の右辺には、第9巻のラプラス変換・たたみこみ（9-6・9-7）を使います」と予告しました。インパルス応答 $g(t)$ がわかっている機械に、時間とともに変わる入力 $u(t)$ を入れたときの出力 $y(t)$ を考えます。

\kixghead{入力を細い柱に分けて、たす}
\begin{enumerate}[label={\bkcircled[\bkTT]{\arabic*}}, leftmargin=2.4\zw, itemsep=2pt, topsep=2pt]
\item 入力を、幅 $\Delta t_1$ の細い柱に分ける。時刻 $t_1$ の柱の面積は $u(t_1)\Delta t_1$ で、これは「時刻 $t_1$ に、力積 $u(t_1)\Delta t_1$ でたたく」のとほぼ同じ（⑦）。
\item 1本の柱に対する出力は、インパルス応答を $t_1$ だけおくらせて、$u(t_1)\Delta t_1$ 倍したもの：$g(t-t_1)\,u(t_1)\Delta t_1$。
\item 機械が線形（5-18）なら、全部の柱の出力をたしたものが、本当の出力。柱を細くした先は積分になる。
\end{enumerate}
\kixgframe[112mm]{\Large $\displaystyle y(t)=\int_0^tg(t-t_1)\,u(t_1)\,\dd t_1=(g*u)(t)$}
これを $g$ と $u$ の\textbf{たたみこみ}といい、$g*u$ と書きます。\bkidx{たたみこみ}{たたみこみ}積分の中の文字 $t_1$ は「前の時刻 $t_1$ に入れた分」です（専門書では $\tau$ と書くことが多いが、この本の $\tau$ は時定数）。

\begin{bkfigure}
\begin{tikzpicture}
\begin{axis}[bkplot, width=112mm, height=50mm, xmin=0, xmax=4, ymin=0, ymax=1.1, xtick={0.5,1,...,4}, ytick={0.5,1}, xlabel={$t$}, ylabel={$y$}]
  \foreach \kixgk in {0,0.5,1,1.5,2,2.5,3,3.5} {
    \edef\kixgtmp{\noexpand\addplot[kixg part, domain=\kixgk:4, samples=40] {0.5*exp(-(x-\kixgk))};}\kixgtmp}
  \addplot[kixg sig, domain=0:4] {1-exp(-x)};
  \node[font=\sffamily\footnotesize, anchor=west, fill=white, inner sep=1pt] at (axis cs:1.9,0.62) {太い線：たたみこみ $1-e^{-t}$};
  \node[font=\sffamily\footnotesize, anchor=west, fill=white, inner sep=1pt] at (axis cs:1.9,0.36) {細い破線：柱1本ずつの出力};
\end{axis}
\end{tikzpicture}
\bkfigcaption{入力 $u=1$（スイッチを入れたまま）を幅0.5の柱に分け、1本ずつの出力 $0.5e^{-(t-t_1)}$（細い破線）を作る。柱を細くして、全部たした先が、太い線 $1-e^{-t}$。}
\label{bk:9-6:fig-tatami}
\end{bkfigure}

\kixghead{$s$ の世界では、かけ算}
たたみこみをラプラス変換すると、かけ算になります（理由は、すぐあとのせいかくに言うと）。
\kixgframe[100mm]{\Large $\mathcal L[g*u]=G(s)\,U(s)$}
図\ref{bk:9-6:fig-tatami}の例 $g=e^{-t}$、$u=1$ で、2通りに確かめます。
\begin{bkhenkei}[0.52]
\int_0^te^{-(t-t_1)}\cdot1\,\dd t_1=e^{-t}\Bigl[e^{t_1}\Bigr]_0^t=1-e^{-t} & 時間の世界：たたみこみの積分 \\
\frac{1}{s+1}\times\frac{1}{s}=\frac{1}{s(s+1)} & $s$ の世界：$G(s)U(s)$ \\
\mathcal L[1-e^{-t}]=\frac{1}{s}-\frac{1}{s+1}=\frac{1}{s(s+1)} & 時間の世界の答えを変換（通分） \\
\end{bkhenkei}
同じになりました。時間の世界のめんどうな積分が、$s$ の世界ではかけ算1回です。これが、9-8 の伝達関数（出てくるもの＝箱×入れるもの）の土台です。なお、ふつうのかけ算 $g\,u$ の変換は $G\,U$ ではありません（9-5 の⑨）。$s$ の世界でかけ算になるのは、たたみこみのほうです。

\begin{seikaku}
\kixgbox{なぜ、たたみこみが、かけ算になるのか}
$\mathcal L[g*u]=\displaystyle\int_0^{\infty}\Bigl(\int_0^tg(t-t_1)u(t_1)\,\dd t_1\Bigr)e^{-st}\,\dd t$ は、$t_1$ と $t$ の平面の三角形 $0\leqq t_1\leqq t$ の上の重積分です。8-6 の場合B のように、積分の順番を入れかえ、先に $t$ で（$t_1$ から $\infty$ まで）積分します。$t-t_1=v$ と置くと $e^{-st}=e^{-st_1}e^{-sv}$ なので、
\[ \int_0^{\infty}u(t_1)e^{-st_1}\Bigl(\int_0^{\infty}g(v)e^{-sv}\,\dd v\Bigr)\dd t_1=G(s)\,U(s) \]
となり、2つの変換のかけ算に分かれます（8-6 の場合 A$'$ と同じ形）。
\end{seikaku}

\bksubsection{9}{初期値の定理・最終値の定理：もどさずに、はじめと最後を知る}

\label{bk:9-6:saishu}$s$ の式を時間の世界にもどすのは 9-7 で学びますが、はじめの値と、最後に落ち着く値だけなら、もどさずにわかります。9-5 の微分の行で、$s$ を0に近づけます。
\begin{bkhenkei}[0.56]
\int_0^{\infty}f'(t)\,e^{-st}\,\dd t=sF(s)-f(0) & 微分の行（9-5） \\
\int_0^{\infty}f'(t)\,\dd t=\lim_{s\to0}sF(s)-f(0) & $s\to0$ で $e^{-st}\to1$ \\
f(\infty)-f(0)=\lim_{s\to0}sF(s)-f(0) & 左辺は、変化をたすと、最後－はじめ（5-11） \\
\end{bkhenkei}
両辺の $f(0)$ が消えて、\textbf{最終値の定理}が出ます。\bkidx{さいしゅうちのていり}{最終値の定理}$s$ を大きくすると $e^{-st}$ で左辺の積分は0に近づくので、\textbf{初期値の定理}も出ます。
\kixgframe[120mm]{\large 最終値の定理\quad $\displaystyle\lim_{t\to\infty}f(t)=\lim_{s\to0}sF(s)$（落ち着くときだけ）\\[4pt]初期値の定理\quad $\displaystyle f(0)=\lim_{s\to\infty}sF(s)$}
\begin{itemize}[leftmargin=1.2\zw, itemsep=2pt, topsep=2pt, label={\color{\bkTT}●}]
\item お湯 $T(s)=\frac{20}{s}+\frac{60}{s+0.05}$：$sT(s)=20+\frac{60s}{s+0.05}$。$s\to0$ で20（最後は部屋の温度）、$s\to\infty$ で $20+60=80$（はじめの温度）。
\item $F(s)=\dfrac{1}{s(s+1)(s+2)}$：$sF(s)=\dfrac{1}{(s+1)(s+2)}\to\dfrac{1}{2}$。9-7 でもどすと $f=\frac12-e^{-t}+\frac12e^{-2t}$ で、$t=5$ で 0.4933、$t=10$ で 0.49995 と、$\frac12$ に近づきます。
\end{itemize}
\textbf{落ち着かないときは使えません}。$\sin t$ の変換 $\frac{1}{s^2+1}$ では $s\cdot\frac{1}{s^2+1}\to0$ ですが、$\sin t$ はずっとゆれて0に落ち着きません。$e^{t}$ でも $\frac{s}{s-1}\to0$ ですが、$e^t$ は限りなく大きくなります。使ってよいのは、$sF(s)$ の分母が0になる点が、すべて $s$ の平面の左半分（実部が負）にあるときです（9-9 の安定）。


\bksubsection{10}{場合分けの表9-6と流れ図}

\label{bk:9-6:baai}ここまでで、変換の求め方がそろいました。9-5 で表9-6のはじめの3行（A・B・E）を使いました。この章で、残りの行（C・D・F〜I）を足して、表を完成させます。I（たたみこみ）は、9-5 の予告にはなかった行です。見開きの左のページが場合分けの表、右のページが流れ図です。流れ図の使い方は、そのあとにあります。

\kixgdanhead{3}{直す前に、楽な道をさがす}
表を引く前に、流れ図の上のほうの問いで、分けたり数値にしたりできないかを見ます。データだけなら数値で（H）、区間ごとに式が変わるならステップ関数の和に分けて（F）、いくつかの項の和なら項ごとに分けて（B）から、表を引きます。変換したあとは、初期値・最終値の定理で、はじめと最後の値をもどさずに確かめます。


\kixgcasetitle{9-6}{ラプラス変換：どんなときに、どう求める？}
\begin{bkcenter}
\begin{bktable}{colspec={Q[c,m,9mm] Q[l,m,46mm] X[l,m] Q[c,m,11mm]}, row{2-Z}={font=\small}, rowsep=2pt, colsep=3pt}
記号 & こんなとき & 手順 & 例題 \\
\kixgcase{A} & 表にある形（$1$、$t^n$、$e^{at}$、$\sin\omega t$、$\cos\omega t$、$\delta(t)$） & \kixgop{2}\enspace 表を引く。まだない行は、定義の積分（\kixgop{1}）で1回だけ作る & 1〜4・7・10・11 \\
\kixgcase{B} & いくつかの項の和（定数倍は前に出す） & \kixgdan{3}\enspace 線形：項ごとに下の問いで変換して、たす & 2・7・11 \\
\kixgcase{C} & $e^{-at}$ に、表にある別の関数（$\sin$・$\cos$・$t^n$ など）がかかっている（$e^{-at}$ だけなら A） & \kixgop{2}\enspace $e^{-at}$ をとった形の $F(s)$ を引き、$s$ を $s+a$ に（$s$ のずらし） & 3 \\
\kixgcase{D} & おくれて始まる（$H(t-t_0)$ がかかっている。$t_0>0$） & \kixgop{2}\enspace おくれをとった形の $F(s)$ を引き、$e^{-st_0}$ をかける（時間のずらし） & 4・7 \\
\kixgcase{E} & $f'$・$f''$ のある式（微分方程式）を変換する & \kixgop{2}\enspace $f'\to sF-f(0)$、$f''\to s^2F-sf(0)-f'(0)$、$\int_0^tf\to\frac{F}{s}$。ほかの項は表で & 5・6 \\
\kixgcase{F} & 区間ごとに式が変わる（パルス・階段） & \kixgdan{3}\enspace ステップ関数 $H(t-t_0)$ の和に直して、B・D へ & 7 \\
\kixgcase{G} & $t^n$ 以外の関数（$\sin$・$\cos$ など）に $t$ がかかっている（$t^n$ だけなら A） & \kixgop{2}\enspace $t$ をとった形の $F(s)$ を引き、$-\dfrac{\dd F}{\dd s}$ & 8 \\
\kixgcase{H} & 表にない形・データだけで、数値で求めたい & \kixgdan{3}\enspace 定義の積分を、表計算の台形で（$e^{-st}$ が十分小さくなる $t$ まで） & 9 \\
\kixgcase{I} & たたみこみ $(f*g)(t)$ を変換する & \kixgop{2}\enspace それぞれを変換して、かけ算 $F(s)\,G(s)$ & 10 \\
\end{bktable}
\end{bkcenter}
{\small A〜D・G・I は直す段（表を引く・ずらす・かけ算）、E は微分方程式を翻訳する所、F・B・H は楽をする段（分ける・数値）です。}

\begin{bkfigure}
\begin{tikzpicture}[x=1mm, y=0.9mm]
  \node[kixg do, text width=80mm] (s) at (46,0) {記号を読む（名前の段）：何を変換するか（$t$ の関数か、\\微分方程式か）。定数倍は前に出す（線形）};
  \node[kixg q] (q1) at (33,-22) {$f'$・$f''$ のある式を変換？};
  \node[kixg do, text width=40mm, anchor=west] (rE) at (71,-22) {$sF-f(0)$ などに。\\ほかの項は表で\enspace\kixgcase{E}};
  \node[kixg q] (q2) at (33,-40) {データだけ・数値で？};
  \node[kixg do, text width=40mm, anchor=west] (rH) at (71,-40) {台形で数値の積分\enspace\kixgcase{H}};
  \node[kixg q] (q3) at (33,-58) {たたみこみの形？};
  \node[kixg do, text width=40mm, anchor=west] (rI) at (71,-58) {それぞれ下の問いで\\変換して、かけ算\enspace\kixgcase{I}};
  \node[kixg q] (q4) at (33,-78) {区間ごとに式が変わる？};
  \node[kixg do, text width=40mm, anchor=west] (rF) at (71,-78) {ステップ関数の和に\\直して、下の問いへ\enspace\kixgcase{F}};
  \node[kixg q] (q5) at (33,-98) {いくつかの項の和？};
  \node[kixg do, text width=40mm, anchor=west] (rB) at (71,-98) {項ごとに下の問いへ。\\あとでたす\enspace\kixgcase{B}};
  \node[kixg q] (q6) at (33,-118) {おくれて始まる？};
  \node[kixg do, text width=40mm, anchor=west] (rD) at (71,-118) {$e^{-st_0}$ をかける。\\おくれをとって下へ\enspace\kixgcase{D}};
  \node[kixg q] (q7) at (33,-138) {$e^{-at}$ に別の関数がかかる？};
  \node[kixg do, text width=40mm, anchor=west] (rC) at (71,-138) {$s$ を $s+a$ に。\\$e^{-at}$ をとって下へ\enspace\kixgcase{C}};
  \node[kixg q] (q8) at (33,-158) {$t^n$ 以外に $t$ がかかる？};
  \node[kixg do, text width=40mm, anchor=west] (rG) at (71,-158) {$t$ をとって表を引き、\\$-\frac{\dd F}{\dd s}$\enspace\kixgcase{G}};
  \node[kixg do, text width=40mm] (rA) at (33,-177) {表を引く\enspace\kixgcase{A}};
  \node[kixg do, text width=110mm] (end) at (62,-196) {確かめ：$s=2$ などで表計算の値と同じか／初期値の定理（$t=0$）・最終値の定理（落ち着くときだけ）};
  \draw[kixg go] (s.south -| q1) -- (q1);
  \foreach \kixgq/\kixgr in {q1/rE,q2/rH,q3/rI,q4/rF,q5/rB,q6/rD,q7/rC,q8/rG} {\draw[kixg go] (\kixgq) -- node[kixg yn, above] {はい} (\kixgr);}
  \foreach \kixgq/\kixgr in {q1/q2,q2/q3,q3/q4,q4/q5,q5/q6,q6/q7,q7/q8,q8/rA} {\draw[kixg go] (\kixgq) -- node[kixg yn, left=1mm] {いいえ} (\kixgr);}
  \draw[kixg go, shorten >=0.4mm] (rI.south) |- (33.8,-69.5);
  \draw[kixg go, shorten >=0.4mm] (rF.south) |- (33.8,-89.5);
  \draw[kixg go, shorten >=0.4mm] (rB.south) |- (33.8,-109.5);
  \draw[kixg go, shorten >=0.4mm] (rD.south) |- (33.8,-129.5);
  \draw[kixg go, shorten >=0.4mm] (rC.south) |- (33.8,-149.5);
  \foreach \kixgn in {rE,rH,rG,rA} {\draw[bkInk!60, line width=0.7pt, shorten <=1.5mm] (\kixgn.east) -- (\kixgn.east -| 127,0);}
  \draw[bkInk!60, line width=0.7pt] (127,-22) -- (127,-196);
  \draw[kixg go, shorten <=0pt] (127,-196) -- (end.east -| 127,0) -- (end.east);
\end{tikzpicture}
\bkfigcaption{ラプラス変換の流れ図。ひし形は問い、四角はすること。I・F・B・D・C のあとは、左の列の次の問いへもどる（B なら項ごとに、D・C ならおくれや $e^{-at}$ をとった形で）。計算した四角からは、右の線をたどって最後の「確かめ」へ。}
\label{bk:9-6:fig-flow}
\end{bkfigure}
\kixghead{流れ図の使い方：$e^{-t}\sin2t$}
\begin{enumerate}[label={\bkcircled[\bkTT]{\arabic*}}, leftmargin=2.4\zw, itemsep=1pt, topsep=2pt]
\item 名前の段：何を＝$e^{-t}\sin2t$（$t$ の関数）。
\item $f'$・$f''$ のある式？\enspace いいえ。\enspace データだけ？\enspace いいえ。\enspace たたみこみ？\enspace いいえ。\enspace 区間ごと？\enspace いいえ。\enspace 項の和？\enspace いいえ。\enspace おくれ？\enspace いいえ。
\item $e^{-at}$ に別の関数がかかる？\enspace はい（$a=1$、別の関数は $\sin2t$）→ \kixgcase{C}。$e^{-t}$ をとった $\sin2t$ で下の問いへ。
\item $t^n$ 以外に $t$ がかかる？\enspace いいえ → \kixgcase{A}：$\sin2t$ の行は $\frac{2}{s^2+4}$。
\item C にもどって、$s$ を $s+1$ にする。答えは次のとおり。
\end{enumerate}
\kixgline{$\mathcal L[e^{-t}\sin2t]=\dfrac{2}{(s+1)^2+4}$}
流れ図の問いの順番にも、わけがあります。微分方程式の問い（E）をいちばん上に置いたのは、微分方程式は「いくつかの項の和」にも見えるので、先に見分けておくためです。また、$e^{-t}$ だけ、$t^3$ だけのように表の1行そのものの関数は、C・G の問いに「いいえ」と答えて A に着きます（$e^{-t}$ を「$e^{-t}\times1$」と見ない）。

\kixghead{どの場合かな？}
流れ図を上からたどって、先に自分で考えてから答え（右）を見ましょう。
\begin{enumerate}[label={\bkcircled[\bkTT]{\arabic*}}, leftmargin=2.4\zw, itemsep=2pt, topsep=2pt]
\item $\mathcal L[\cos5t]$\hfill 答え：\kixgcase{A}（表の $\cos\omega t$ の行）
\item $\mathcal L[e^{-3t}\cos t]$\hfill 答え：\kixgcase{C}→\kixgcase{A}
\item 3秒後に 12\,V をつなぐ電圧 $12H(t-3)$\hfill 答え：\kixgcase{D}→\kixgcase{A}
\item はじめの1秒間だけ 5\,V をかける電圧\hfill 答え：\kixgcase{F}→\kixgcase{B}→（\kixgcase{A} と \kixgcase{D}→\kixgcase{A}）
\item $x'+4x=0$、$x(0)=3$ の両辺の変換\hfill 答え：\kixgcase{E}
\end{enumerate}
$\mathcal L[t^2e^{-t}]$ なら、$e^{-t}$ に $t^2$ がかかっているので C に行き、$e^{-t}$ をとると $t^2$ だけが残ります。$t^2$ は「$t^n$ 以外に $t$ がかかる」ではないので、A（$t^n$ の行）に着きます。答えは $\frac{2}{s^3}$ の $s$ を $s+1$ にした $\frac{2}{(s+1)^3}$ です。



\kixghead{Python で、表の行を確かめる}
表の行が正しいかは、$s$ に数を入れて、定義の積分を数値で計算すればわかります（9-5 と同じ \texttt{quad}）。$s=2$ で、$t^3$、$e^{-t}\sin2t$、$t\sin2t$、$e^{-(t-1)}H(t-1)$ の4行を確かめます。$\infty$ のかわりに $t=50$ まで積分します（$e^{-100}$ はほとんど0）。
\Needspace{17\baselineskip}
\begin{bkcode}[Python]
import numpy as np
from scipy.integrate import quad
s = 2
rows = [(lambda t: t**3, 6/s**4),
        (lambda t: np.exp(-t)*np.sin(2*t), 2/((s+1)**2+4)),
        (lambda t: t*np.sin(2*t), 4*s/(s**2+4)**2),
        (lambda t: (t >= 1)*np.exp(-(t-1)), np.exp(-s)/(s+1))]
for f, F in rows:                \# f は関数、F は表の値
    v = quad(lambda t: f(t)*np.exp(-s*t), 0, 50, limit=200)[0]
    print(round(v, 5), round(F, 5))
\end{bkcode}
\begin{bkcodeout}[実行結果]
0.375 0.375
0.15385 0.15385
0.125 0.125
0.04511 0.04511
\end{bkcodeout}
左が数値の積分、右が表の式の値です。4つの行とも同じになりました。\texttt{(t >= 1)} は、$t\geqq1$ のとき1、そうでないとき0になる書き方で、$H(t-1)$ のかわりです。


\begin{motto}[5]
\kixgbox{電気の回路で見る、たたみこみ}
RC 回路（$\tau=RC$）のコンデンサーの電圧は、電池の電圧 $e(t)$ が時間とともに変わっても、インパルス応答 $g(t)=\frac{1}{\tau}e^{-t/\tau}$ とのたたみこみ $v=g*e$ で求められます。「少し前に入れた電圧ほど、効き目が $e^{-(t-t_1)/\tau}$ で小さくなって残っている」と読めます。コンデンサーが、過去の入力を少しずつ忘れながら覚えている、ということです。6-13 で「$\cos$ でない一般の形の右辺」と言ったのは、このような入力のことです。
\end{motto}

\bksubsection{11}{変換表（見開き）}

\label{bk:9-6:hyou}9-5 とこの章で作った行の表です。左は関数、右は性質の表です。9-7 からは、この表を引いて計算します。\bkhatenaref の答えは、表の1・3・5・7〜10行めです：定数 $C\to\frac{C}{s}$、$t\to\frac{1}{s^2}$、$e^{at}\to\frac{1}{s-a}$、$\sin\omega t\to\frac{\omega}{s^2+\omega^2}$、$\cos\omega t\to\frac{s}{s^2+\omega^2}$、$e^{-at}\sin\omega t\to\frac{\omega}{(s+a)^2+\omega^2}$。

\begin{bkcenter}
\begin{bktable}{colspec={Q[c,m,30mm] Q[c,m,34mm] X[l,m] Q[c,m,15mm]}, row{2-Z}={font=\small}, rowsep=1.6pt, colsep=3pt}
$f(t)$ & $F(s)=\mathcal L[f]$ & 作り方 & $s$ の条件 \\
$\delta(t)$ & $1$ & 細いパルスの先（9-6⑦） & どの $s$ も \\
$1$（$H(t)$） & $\dfrac{1}{s}$ & 定義の積分（9-5） & $s>0$ \\
$t$ & $\dfrac{1}{s^2}$ & 部分積分（9-5） & $s>0$ \\
$t^n$ & $\dfrac{n!}{s^{n+1}}$ & 部分積分のくり返し（9-6②） & $s>0$ \\
$e^{at}$ & $\dfrac{1}{s-a}$ & 定義の積分（9-5） & $s>a$ \\
$t\,e^{at}$ & $\dfrac{1}{(s-a)^2}$ & $s$ のずらし（9-6④） & $s>a$ \\
$\sin\omega t$ & $\dfrac{\omega}{s^2+\omega^2}$ & オイラーの公式（9-6③） & $s>0$ \\
$\cos\omega t$ & $\dfrac{s}{s^2+\omega^2}$ & オイラーの公式（9-6③） & $s>0$ \\
$e^{-at}\sin\omega t$ & $\dfrac{\omega}{(s+a)^2+\omega^2}$ & $s$ のずらし（9-6④） & $s>-a$ \\
$e^{-at}\cos\omega t$ & $\dfrac{s+a}{(s+a)^2+\omega^2}$ & $s$ のずらし（9-6④） & $s>-a$ \\
$t\sin\omega t$ & $\dfrac{2\omega s}{(s^2+\omega^2)^2}$ & $t$ をかける行（9-6⑥） & $s>0$ \\
$t\cos\omega t$ & $\dfrac{s^2-\omega^2}{(s^2+\omega^2)^2}$ & $t$ をかける行（9-6⑥） & $s>0$ \\
$H(t-t_0)$ & $\dfrac{e^{-st_0}}{s}$ & 時間のずらし（9-6⑤） & $s>0$ \\
パルス（高さ1、0〜$t_0$） & $\dfrac{1-e^{-st_0}}{s}$ & ステップの和（9-6⑤） & どの $s$ も \\
\end{bktable}
\bktabcaption{ラプラス変換の表（関数）。$a$、$\omega$、$t_0$ は定数（$\omega>0$、$t_0>0$）。$s$ の条件は、$s$ が複素数なら「実部が」と読む。}
\label{bk:9-6:tab-kansu}
\end{bkcenter}

\begin{bkcenter}
\begin{bktable}{colspec={Q[l,m,28mm] Q[c,m,34mm] Q[c,m,32mm] X[l,m]}, row{2-Z}={font=\small}, rowsep=4pt, colsep=3pt}
性質 & 時間の世界 & $s$ の世界 & 作り方 \\
線形 & $af+bg$ & $aF+bG$ & 積分の性質（9-5⑨） \\
微分 & $f'$ & $sF-f(0)$ & 部分積分（9-5⑥） \\
2回微分 & $f''$ & $s^2F-sf(0)-f'(0)$ & 微分の行を2回（9-5⑦） \\
積分 & $\displaystyle\int_0^tf\,\dd t_1$ & $\dfrac{F}{s}$ & 微分の行から（9-6⑥） \\
$s$ のずらし & $e^{-at}f(t)$ & $F(s+a)$ & 肩のたし算（9-6④） \\
時間のずらし & $f(t-t_0)H(t-t_0)$ & $e^{-st_0}F(s)$ & 置換積分（9-6⑤） \\
$t$ をかける & $t\,f(t)$ & $-\dfrac{\dd F}{\dd s}$ & $s$ で微分（9-6⑥） \\
たたみこみ & $(f*g)(t)$ & $F(s)\,G(s)$ & 積分の順番を入れかえ（9-6⑧） \\
初期値の定理 & $f(0)$ & $\displaystyle\lim_{s\to\infty}sF(s)$ & 微分の行から（9-6⑨） \\
最終値の定理 & $\displaystyle\lim_{t\to\infty}f(t)$ & $\displaystyle\lim_{s\to0}sF(s)$ & 落ち着くときだけ（9-6⑨） \\
\end{bktable}
\bktabcaption{ラプラス変換の表（性質）。ふつうのかけ算 $fg$ の行はない（$FG$ にはならない）。}
\label{bk:9-6:tab-seishitsu}
\end{bkcenter}
\noindent 表を引くときは、表の形と「文字の置き方」まで同じかを確かめます。たとえば $\frac{3}{s+2}$ は、$e^{at}$ の行の3倍で $a=-2$（$3e^{-2t}$）です。$\frac{1}{s^2+9}$ は $\sin3t$ の行 $\frac{3}{s^2+9}$ の $\frac13$ 倍（$\frac13\sin3t$）です。分子の数を合わせてから引きます（9-7 で何度も使います）。

\begin{senmonsho}[変換表の専門書ふうの書き方]
\begin{bkhonbun}
\[ \mathcal{L}\{u(t-a)\}=\frac{e^{-as}}{s},\qquad \mathcal{L}\{\delta(t)\}=1. \]
$f(t)$ と $g(t)$ の畳み込みを
\[ (f*g)(t)=\int_0^t f(\tau)g(t-\tau)\,d\tau \]
で定義すると，$\mathcal{L}\{f*g\}=F(s)G(s)$ が成り立つ．
\end{bkhonbun}
\bkyaku{$u(t-a)$ は、この本のステップ関数 $H(t-t_0)$（$u$ や $1(t)$ と書く本も多い）。「畳み込み」は、この本の「たたみこみ」。積分の中の $\tau$ は、この本の $t_1$（この本の $\tau$ は時定数）。$f(\tau)g(t-\tau)$ と、$g$ のほうをずらして書いても、同じ値になる（$t_1$ を $t-t_1$ に置きかえるとわかる）。}
\end{senmonsho}

\begin{kotoba}
\small
\bkword{変換表}{時間の関数と、そのラプラス変換の表。関数の表と性質の表がある。行は1回だけ作り、あとは引く。}
\bkword{$s$ のずらし}{$e^{-at}$ をかけると、$F(s)$ の $s$ が $s+a$ になる。}
\bkword{時間のずらし}{$t_0$ おくれて始まると、$F(s)$ に $e^{-st_0}$ がかかる。むだ時間は $e^{-sL}$。}
\bkword{デルタ関数}{面積1のまま幅を細くした先。$\delta(t)$。たたいた力積を表す。$\mathcal L[\delta]=1$。}
\bkword{インパルス応答}{$\delta(t)$ でたたいたときの出力。}
\bkword{たたみこみ}{$(g*u)(t)=\int_0^tg(t-t_1)u(t_1)\,\dd t_1$。変換するとかけ算。}
\bkword{最終値の定理}{落ち着くときだけ、$\lim_{t\to\infty}f=\lim_{s\to0}sF$。}
\end{kotoba}

\bkkigocards{\bkkigocard{$\delta(t)$}{デルタ・ティー（デルタ関数）}{面積1で幅が限りなく細いパルス。たたいた力積。6-7 の $\delta$ とは別}{Dirac delta function}{9-6}}
{\bkkigocard{$f*g$}{エフ・たたみこみ・ジー（$*$ をスターと読む人もいる）}{前の時刻の入力の効き目をたし合わせたもの。変換すると $F(s)G(s)$}{convolution}{9-6}}

\begin{taisetsu}
\kixgbox{ラプラス変換の3つの段}
\kixgdan{1}\enspace 記号を読む：何を変換するか（$t$ の関数か、微分方程式か）\par
\kixgdan{3}\enspace 直す前に：データなら数値で、区間ごとならステップの和、項の和なら項ごとに\par
\kixgdan{2}\enspace 表を引く。ずれ・おくれ・$t$ は性質の表で表の形にもどしてから\par
\kixgbox{場合分けの表9-6（まとめ）}
\begin{bktable}{colspec={Q[c,m,9mm] Q[l,m,42mm] X[l,m]}, row{2-Z}={font=\small}, rowsep=1.5pt, colsep=3pt}
記号 & こんなとき & 直し方 \\
\kixgcase{A} & 表にある形 & 表を引く \\
\kixgcase{B} & 項の和 & 項ごとに変換して、たす \\
\kixgcase{C} & $e^{-at}\times f$ & $F(s+a)$ \\
\kixgcase{D} & おくれて始まる & $e^{-st_0}F(s)$ \\
\kixgcase{E} & 微分方程式 & $sF-f(0)$ など \\
\kixgcase{F} & 区間ごと & ステップ関数の和 \\
\kixgcase{G} & $t\times f$ & $-\frac{\dd F}{\dd s}$ \\
\kixgcase{H} & データ・数値 & 台形で数値の積分 \\
\kixgcase{I} & たたみこみ & $F(s)G(s)$ \\
\end{bktable}
\end{taisetsu}

%% ---------------- 例題 ----------------
\Needspace{20\baselineskip}\bksection{reidai}{例題でたしかめよう}

\noindent 例題の「考え方」の1行めには、場合分けの表9-6のどの場合かを書きました。

\kixghead{例題を、流れ図でたどってみよう}
例題1〜11を、図\ref{bk:9-6:fig-flow}の流れ図でたどると、次のようになります（書いていない問いは「いいえ」）。
\begin{bkcenter}
\begin{bktable}{colspec={Q[c,m,10mm] X[l,m] Q[c,m,24mm]}, row{2-Z}={font=\small}, rowsep=0pt, colsep=3pt}
例題 & 「はい」と答えた問い & 着いた所 \\
1 & なし & \kixgcase{A} \\
2 & 項の和？ → 項ごとに A & \kixgcase{B}\kixgcase{A} \\
3 & $e^{-at}$ に別の関数がかかる？ → A & \kixgcase{C}\kixgcase{A} \\
4 & おくれて始まる？ → おくれをとって A & \kixgcase{D}\kixgcase{A} \\
5 & $f'$・$f''$ のある式を変換？ & \kixgcase{E} \\
6 & $f'$・$f''$ のある式を変換？ & \kixgcase{E} \\
7 & 区間ごと？ → 項の和？ → 項ごとに A・D & \kixgcase{F}\kixgcase{B}\kixgcase{D}\kixgcase{A} \\
8 & $t^n$ 以外に $t$ がかかる？ & \kixgcase{G} \\
9 & データだけ・数値で？ & \kixgcase{H} \\
10 & たたみこみ？ → それぞれ A & \kixgcase{I}\kixgcase{A} \\
11 & 項の和？ → 項ごとに A → 確かめ & \kixgcase{B}\kixgcase{A} \\
\end{bktable}
\bktabcaption{例題1〜11の流れ図の道すじ（じぶんでとこう・章末問題も、同じ番号の例題と同じ）。}
\label{bk:9-6:tab-michi}
\end{bkcenter}

\begin{reidai}[表を引く]
次のラプラス変換を、表を引いて求めましょう。（1）$\mathcal L[t^3]$\quad（2）$\mathcal L[\sin3t]$\quad（3）$\mathcal L[\cos3t]$\par
\kangaekata \kixgbaai{表にある形}{9-6}{A}\par
\kixgdan{1}\enspace 名前は $\mathcal L$ だが、ここでは積分しない（行は②・③で作ってある）。\par
\kixgnaosu{2}\enspace $t^n$ の行で $n=3$、$\sin\omega t$・$\cos\omega t$ の行で $\omega=3$。\par
\shiki\par\nobreak
\begin{kixghenkei}[0.5]
\mathcal L[t^3]=\frac{3!}{s^{4}}=\frac{6}{s^4} & $\frac{n!}{s^{n+1}}$ で $n=3$ \\
\mathcal L[\sin3t]=\frac{3}{s^2+3^2}=\frac{3}{s^2+9} & 分子は $\omega=3$ \\
\mathcal L[\cos3t]=\frac{s}{s^2+9} & 分子は $s$ \\
\end{kixghenkei}
\kotae （1）\nobreak $\dfrac{6}{s^4}$　（2）\nobreak $\dfrac{3}{s^2+9}$　（3）\nobreak $\dfrac{s}{s^2+9}$
\end{reidai}

\begin{jibun}
次のラプラス変換を、表を引いて求めましょう。\par
（1）\nobreak $\mathcal L[t^2]$\par
（2）\nobreak $\mathcal L[\sin5t]$\par
（3）\nobreak $\mathcal L[\cos2t]$
\bkjans{（1）\nobreak $\frac{2}{s^3}$。（2）\nobreak $\frac{5}{s^2+25}$。（3）\nobreak $\frac{s}{s^2+4}$。}
\end{jibun}

\begin{reidai}[項の和]
$\mathcal L[2+3t^2-4\cos2t]$ を求めましょう。\par
\kangaekata \kixgbaai{いくつかの項の和}{9-6}{B,A}\par
\kixgdan{1}\enspace 名前は $\mathcal L$。ここでは積分しない。\par
\kixgdan{3}\enspace 線形なので、項ごとに変換して、たす。定数倍は前に出す。\par
\kixgnaosu{2}\enspace $1\to\tfrac{1}{s}$、$t^2\to\tfrac{2}{s^3}$、$\cos2t\to\tfrac{s}{s^2+4}$。\par
\shiki $2\times\dfrac{1}{s}+3\times\dfrac{2}{s^3}-4\times\dfrac{s}{s^2+4}$\par
\kotae $\dfrac{2}{s}+\dfrac{6}{s^3}-\dfrac{4s}{s^2+4}$
\end{reidai}

\begin{reidai}[ゆれながら止まる形]
次のラプラス変換を求めましょう。（1）$\mathcal L[e^{-t}\sin2t]$\quad（2）$\mathcal L[e^{-t}\cos2t]$\par
\kangaekata \kixgbaai{$e^{-at}$ とほかの関数のかけ算}{9-6}{C,A}\par
\kixgdan{1}\enspace 名前は $\mathcal L$。ここでは積分しない。\par
\kixgnaosu{2}\enspace $e^{-t}$ をとった $\sin2t$・$\cos2t$ の行を引き、$s$ を $s+1$ にする。\par
\shiki\par\nobreak
\begin{kixghenkei}[0.52]
\mathcal L[\sin2t]=\frac{2}{s^2+4}\ \to\ \frac{2}{(s+1)^2+4} & （1）$s$ を $s+1$ に \\
\mathcal L[\cos2t]=\frac{s}{s^2+4}\ \to\ \frac{s+1}{(s+1)^2+4} & （2）分子の $s$ も $s+1$ に \\
\end{kixghenkei}
\kotae （1）\nobreak $\dfrac{2}{(s+1)^2+4}=\dfrac{2}{s^2+2s+5}$　（2）\nobreak $\dfrac{s+1}{s^2+2s+5}$（分母は展開した）
\end{reidai}

\begingroup\tcbset{equal height group=kixgJB}%
\noindent\begin{minipage}[t]{0.485\linewidth}\vspace{0pt}
\begin{jibun}
\raggedright $\mathcal L[1-t+\sin t]$ を求めましょう。
\bkjans{$\frac{1}{s}-\frac{1}{s^2}+\frac{1}{s^2+1}$。}
\end{jibun}
\end{minipage}\hfill
\begin{minipage}[t]{0.485\linewidth}\vspace{0pt}
\begin{jibun}
\raggedright 次のラプラス変換を求めましょう。\par
（1）\nobreak $\mathcal L[e^{-2t}\sin3t]$\par
（2）\nobreak $\mathcal L[e^{-2t}\cos3t]$
\bkjans{（1）\nobreak $\frac{3}{(s+2)^2+9}=\frac{3}{s^2+4s+13}$。（2）\nobreak $\frac{s+2}{s^2+4s+13}$。}
\end{jibun}
\end{minipage}\par
\endgroup

\begin{reidai}[おくれて入れる]
次の関数のラプラス変換を求めましょう。（1）2秒後にスイッチを入れて 10\,V をつなぐ電圧 $10H(t-2)$\quad（2）$e^{-(t-1)}H(t-1)$\quad（3）3秒後に力積 \qty{5}{N.s} でたたく力 $5\delta(t-3)$\par
\kangaekata \kixgbaai{おくれて始まる}{9-6}{D,A}\par
\kixgdan{1}\enspace 名前は $\mathcal L$。ここでは積分しない。\par
\kixgnaosu{2}\enspace おくれをとった形（10、$e^{-t}$、$5\delta(t)$）の行を引き、$e^{-st_0}$ をかける。\par
\shiki\par\nobreak
\begin{kixghenkei}[0.5]
\mathcal L[10]=\frac{10}{s}\ \to\ \frac{10e^{-2s}}{s} & （1）$t_0=2$ なので $e^{-2s}$ をかける \\
\mathcal L[e^{-t}]=\frac{1}{s+1}\ \to\ \frac{e^{-s}}{s+1} & （2）$t_0=1$ なので $e^{-s}$ をかける \\
\mathcal L[5\delta(t)]=5\times1\ \to\ 5e^{-3s} & （3）$t_0=3$ なので $e^{-3s}$ をかける \\
\end{kixghenkei}
\kotae （1）\nobreak $\dfrac{10e^{-2s}}{s}$　（2）\nobreak $\dfrac{e^{-s}}{s+1}$　（3）\nobreak $5e^{-3s}$
\end{reidai}

\begin{jibun}
次の関数のラプラス変換を求めましょう。\par
（1）\nobreak $5H(t-3)$\par
（2）\nobreak $\sin(t-2)\,H(t-2)$\par
（3）\nobreak 1秒後に力積 \qty{2}{N.s} でたたく力 $2\delta(t-1)$
\bkjans{（1）\nobreak $\frac{5e^{-3s}}{s}$。（2）\nobreak $\frac{e^{-2s}}{s^2+1}$。（3）\nobreak $2e^{-s}$。}
\end{jibun}

\begin{reidai}[車B の式を変換する]
車B の式とはじめの値は、次のとおりです。両辺をラプラス変換して、$X(s)$ を求めましょう。
\[ \ddot x+2\dot x+5x=0,\qquad x(0)=1,\quad \dot x(0)=0 \]
\kangaekata \kixgbaai{微分方程式を変換する}{9-6}{E}\par
\kixgdan{1}\enspace 両辺の $\mathcal L$ をとる。\par
\kixgnaosu{2}\enspace $\ddot x$ は $s^2X-s\cdot1-0$、$\dot x$ は $sX-1$、$x$ は $X$。\par
\shiki\par\nobreak
\begin{kixghenkei}[0.5]
(s^2X-s)+2(sX-1)+5X=0 & 1項ずつ表で \\
(s^2+2s+5)\,X=s+2 & $X$ でくくり、$-s-2$ を右へ \\
X(s)=\frac{s+2}{s^2+2s+5} & わる \\
=\frac{s+2}{(s+1)^2+4} & 分母は例題3と同じ形 \\
\end{kixghenkei}
確かめ（初期値の定理）：$sX(s)=\dfrac{s^2+2s}{s^2+2s+5}$ は、$s$ を大きくすると1に近づく。$x(0)=1$ と合っています。\par
\kotae $X(s)=\dfrac{s+2}{s^2+2s+5}$（分母は 6-12 の車B の特性方程式の左辺。もどし方は 9-7）
\end{reidai}

\begin{reidai}[インパルス応答]
$x'+2x=\delta(t)$、はじめ $x=0$ の両辺をラプラス変換して $X(s)$ を求め、表を逆に引いて $x$ を求めましょう。\par
\kangaekata \kixgbaai{微分方程式を変換する}{9-6}{E}\par
\kixgdan{1}\enspace 両辺の $\mathcal L$ をとる。\par
\kixgnaosu{2}\enspace $x'$ は $sX-0$、$\delta(t)$ は1。\par
\shiki\par\nobreak
\begin{kixghenkei}[0.5]
sX+2X=1 & 1項ずつ表で \\
X(s)=\frac{1}{s+2} & $s+2$ でわる \\
x=e^{-2t} & 表を逆に引く（$e^{at}$ の行で $a=-2$） \\
\end{kixghenkei}
\begin{center}\begin{tikzpicture}
\begin{axis}[bkplot, width=80mm, height=30mm, xmin=-0.5, xmax=3, ymin=0, ymax=1.25, xtick={1,2,3}, ytick={1}, xlabel={$t$}, ylabel={$x$}]
  \addplot[bkcurve, domain=-0.5:0] {0};
  \addplot[bkcurve, domain=0:2.9, samples=60] {exp(-2*x)};
  \draw[bkBlue, line width=0.8pt, dash pattern=on 1.5pt off 1pt] (axis cs:0,0) -- (axis cs:0,1);
  \node[font=\sffamily\footnotesize, anchor=west] at (axis cs:0.9,0.7) {たたいた瞬間に1、そのあと $e^{-2t}$};
\end{axis}
\end{tikzpicture}\end{center}
\kotae $X(s)=\dfrac{1}{s+2}$、$x=e^{-2t}$
\end{reidai}

\begingroup\tcbset{equal height group=kixgJE}%
\noindent\begin{minipage}[t]{0.485\linewidth}\vspace{0pt}
\begin{jibun}
\raggedright $\ddot x+4\dot x+13x=0$、はじめ $x(0)=1$、$\dot x(0)=0$ の両辺をラプラス変換して、$X(s)$ を求めましょう。
\bkjans{$(s^2X-s)+4(sX-1)+13X=0$。$X(s)=\frac{s+4}{s^2+4s+13}=\frac{s+4}{(s+2)^2+9}$。}
\end{jibun}
\end{minipage}\hfill
\begin{minipage}[t]{0.485\linewidth}\vspace{0pt}
\begin{jibun}
\raggedright $x'+0.5x=4\delta(t)$、はじめ $x=0$ の両辺をラプラス変換して $X(s)$ を求め、表を逆に引いて $x$ を求めましょう。
\bkjans{$sX+0.5X=4$。$X(s)=\frac{4}{s+0.5}$。$x=4e^{-0.5t}$。}
\end{jibun}
\end{minipage}\par
\endgroup

\begin{reidai}[パルスを変換する]
はじめの2秒間だけ 3\,V をかけ、そのあと 0\,V にする電圧のラプラス変換を求めましょう。また、$s=2$ のときの値を求めましょう（$e^{-4}=0.0183$）。\par
\kangaekata \kixgbaai{区間ごとに式が変わる}{9-6}{F,B,D,A}\par
\kixgdan{1}\enspace 名前は $\mathcal L$。ここでは積分しない。\par
\kixgdan{3}\enspace 「時刻0に3を足し、時刻2に3をひく」：$3H(t)-3H(t-2)$。項ごとに。\par
\kixgnaosu{2}\enspace $H(t)\to\tfrac{1}{s}$（A）、$H(t-2)\to\tfrac{e^{-2s}}{s}$（D）。\par
\shiki\par\nobreak
\begin{kixghenkei}[0.5]
\mathcal L[3H(t)-3H(t-2)]=\frac{3}{s}-\frac{3e^{-2s}}{s} & 項ごとに \\
=\frac{3(1-e^{-2s})}{s} & まとめる \\
s=2:\ \frac{3(1-0.0183)}{2}=1.4725 & $e^{-2\times2}=e^{-4}$ \\
\end{kixghenkei}
\begin{center}\begin{tikzpicture}
\begin{axis}[bkplot, width=80mm, height=28mm, xmin=-0.5, xmax=4, ymin=0, ymax=3.8, xtick={1,2,3,4}, ytick={3}, xlabel={$t$（s）}, ylabel={電圧（V）}]
  \addplot[bkcurve, domain=-0.5:0] {0};
  \addplot[bkcurve, domain=0:2] {3};
  \addplot[bkcurve, domain=2:3.9] {0};
  \draw[bkBlue, line width=0.8pt, dash pattern=on 1.5pt off 1pt] (axis cs:0,0) -- (axis cs:0,3) (axis cs:2,0) -- (axis cs:2,3);
\end{axis}
\end{tikzpicture}\end{center}
途中から始まるパルス（たとえば時刻1秒から3秒までだけ 5\,V）なら、「時刻1に5を足し、時刻3に5をひく」で $5H(t-1)-5H(t-3)$ と、おくれたステップ2つの和になり、どちらも D で変換します。\par
\kotae $\dfrac{3(1-e^{-2s})}{s}$、$s=2$ で約1.47
\end{reidai}

\begin{jibun}
時刻1秒から3秒までの間だけ 5\,V をかけ、ほかの時刻は 0\,V の電圧のラプラス変換を求めましょう。
\bkjans{$5H(t-1)-5H(t-3)$。$\frac{5e^{-s}}{s}-\frac{5e^{-3s}}{s}=\frac{5(e^{-s}-e^{-3s})}{s}$。}
\end{jibun}

\begin{reidai}[$t$ がかかった形]
$\mathcal L[t\sin2t]$ を求めましょう。\par
\kangaekata \kixgbaai{$t$ とほかの関数のかけ算}{9-6}{G}\par
\kixgdan{1}\enspace 名前は $\mathcal L$。ここでは積分しない。\par
\kixgnaosu{2}\enspace $t$ をとった $\sin2t$ の行 $F=\tfrac{2}{s^2+4}$ を引き、$-\tfrac{\dd F}{\dd s}$。\par
\shiki\par\nobreak
\begin{kixghenkei}[0.5]
F(s)=2(s^2+4)^{-1} & 微分しやすい形に \\
\frac{\dd F}{\dd s}=2\times(-1)(s^2+4)^{-2}\times2s & 合成関数の微分（5-6） \\
-\frac{\dd F}{\dd s}=\frac{4s}{(s^2+4)^2} & $-1$ をかける \\
\end{kixghenkei}
\begin{center}\begin{tikzpicture}
\begin{axis}[bkplot, axis lines=left, width=84mm, height=44mm, xmin=0, xmax=9.5, ymin=-10, ymax=13, xtick={0,2,4,6,8}, ytick={-5,0,5}, xlabel={$t$}]
  \addplot[bkthin, dashed, domain=0:9.2, samples=40] {x};
  \addplot[bkthin, dashed, domain=0:9.2, samples=40] {-x};
  \addplot[bkcurve, domain=0:9.2, samples=240] {x*sin(deg(2*x))};
  \node[font=\sffamily\footnotesize, anchor=west, fill=white, inner sep=1pt] at (axis cs:0.4,11.5) {$t\sin2t$：$t$ に比例して大きくなるゆれ};
\end{axis}
\end{tikzpicture}\end{center}
\kotae $\dfrac{4s}{(s^2+4)^2}$（$s=2$ で $\frac{8}{64}=0.125$。数値の積分も 0.125）。6-13 の共振 $\frac{t}{4}\sin2t$ もこの形（9-7）
\end{reidai}

\begin{jibun}
$\mathcal L[t\sin t]$ を求めましょう。
\bkjans{$F=\frac{1}{s^2+1}=(s^2+1)^{-1}$。$\frac{\dd F}{\dd s}=-(s^2+1)^{-2}\times2s$。$-\frac{\dd F}{\dd s}=\frac{2s}{(s^2+1)^2}$。}
\end{jibun}

\begin{reidai}[表にない形を、数値で]
$f(t)=\dfrac{1}{1+t}$ のラプラス変換の、$s=1$ のときの値を、表計算の台形（$\Delta t=0.1$、$t=10$ まで）で、数値で求めましょう。\par
\kangaekata \kixgbaai{表にない形を、数値で}{9-6}{H}\par
\kixgdan{1}\enspace 名前は $\mathcal L$：$\frac{1}{1+t}$ に $e^{-t}$ をかけて、0から $\infty$ までたす。\par
\kixgdan{3}\enspace 表にない形なので、数値で。$t=10$ で $e^{-10}=0.000045$ なので、そこで止める。\par
\shiki\par\nobreak
\begin{bkcenter}
\bksheet{4}{20}{%
1 & \SetCell{c} $t$ & \SetCell{c} $\frac{1}{1+t}$ & \SetCell{c} $e^{-t}$ & \SetCell{c} かけた値 \\
2 & 0 & 1.0000 & 1.0000 & 1.0000 \\
3 & 0.1 & 0.9091 & 0.9048 & 0.8226 \\
4 & 0.2 & 0.8333 & 0.8187 & 0.6823 \\
$\vdots$ & $\vdots$ & $\vdots$ & $\vdots$ & $\vdots$ \\
102 & 10 & 0.0909 & 0.0000 & 0.0000 \\
103 & \SetCell{c} 台形 & & & \SetCell{bg=kixgS} 0.5980 \\}
\end{bkcenter}
D103 は \texttt{=0.1*(SUM(D2:D102)-(D2+D102)/2)}（台形の和：両はしは半分）。\par

\kotae 約0.598（$\Delta t=0.01$ にすると 0.5964 になり、本当の値 0.59635 に近づく）
\end{reidai}

\begin{jibun}
$f(t)=\dfrac{1}{(1+t)^2}$ のラプラス変換の、$s=1$ のときの値を、表計算の台形（$\Delta t=0.1$、$t=10$ まで）で求めましょう。
\bkjans{かけた値は $t=0$ で1、0.1 で0.7478、0.2 で0.5686…。台形の和は約0.406（本当の値 0.4037）。}
\end{jibun}

\begin{reidai}[たたみこみを2通りで]
$e^{-t}*e^{-2t}$ を、（1）たたみこみの積分で、（2）$s$ の世界のかけ算で求め、同じになることを確かめましょう。\par
\kangaekata \kixgbaai{たたみこみ}{9-6}{I,A}\par
\kixgdan{1}\enspace $(e^{-t}*e^{-2t})(t)=\int_0^te^{-(t-t_1)}e^{-2t_1}\,\dd t_1$。$\mathcal L$ はかけ算になる。\par
\kixgnaosu{1}\enspace （1）は $t_1$ の定積分。$e^{-t}$ は $t_1$ に関係しないので外へ。\par
\shiki\par\nobreak
\begin{kixghenkei}[0.54]
\int_0^te^{-t}e^{t_1}e^{-2t_1}\,\dd t_1=e^{-t}\int_0^te^{-t_1}\,\dd t_1 & （1）肩を分けて、$e^{-t}$ を外へ \\
=e^{-t}\bigl(1-e^{-t}\bigr)=e^{-t}-e^{-2t} & $\int_0^te^{-t_1}\dd t_1=1-e^{-t}$ \\
\frac{1}{s+1}\times\frac{1}{s+2}=\frac{1}{(s+1)(s+2)} & （2）それぞれ表で（A）、かけ算 \\
\mathcal L[e^{-t}-e^{-2t}]=\frac{(s+2)-(s+1)}{(s+1)(s+2)} & （1）の答えを変換して通分 \\
\end{kixghenkei}
\kotae $e^{-t}-e^{-2t}$。変換は $\dfrac{1}{(s+1)(s+2)}$ で、（2）と同じ
\end{reidai}

\begin{jibun}
$1*e^{-3t}$ を、たたみこみの積分と、$s$ の世界のかけ算で求め、同じになることを確かめましょう。
\bkjans{積分：$\int_0^te^{-3t_1}\cdot1\,\dd t_1=\frac{1-e^{-3t}}{3}$（$g=1$ を前の時刻にずらしても1）。かけ算：$\frac{1}{s}\cdot\frac{1}{s+3}$。$\mathcal L[\frac{1-e^{-3t}}{3}]=\frac13(\frac1s-\frac{1}{s+3})=\frac{1}{s(s+3)}$。同じ。}
\end{jibun}

\begin{reidai}[はじめと最後の値で確かめる]
$f(t)=1-e^{-2t}$ を変換して $F(s)$ を求め、初期値の定理と最終値の定理で、$t=0$ の値と最後の値を確かめましょう。\par
\kangaekata \kixgbaai{いくつかの項の和}{9-6}{B,A}\par
\kixgdan{1}\enspace 名前は $\mathcal L$。ここでは積分しない。\par
\kixgdan{3}\enspace 項ごとに表を引く。確かめは、もどさずに $sF(s)$ の $s\to\infty$ と $s\to0$。\par
\shiki\par\nobreak
\begin{kixghenkei}[0.5]
F(s)=\frac{1}{s}-\frac{1}{s+2}=\frac{2}{s(s+2)} & 項ごとに表で。通分 \\
sF(s)=\frac{2}{s+2} & $s$ をかける \\
s\to\infty\text{ で }\frac{2}{s+2}\to0 & 初期値の定理：はじめは0 \\
s\to0\text{ で }\frac{2}{s+2}\to1 & 最終値の定理（$-2$ は左半分） \\
\end{kixghenkei}
\kotae $F(s)=\dfrac{2}{s(s+2)}$。はじめ0、最後は1で、$1-e^{-2t}$ の $t=0$ の値・最後の値と合う
\end{reidai}

\begin{jibun}
$f(t)=10-10e^{-0.5t}$（RC の充電）を変換して $F(s)$ を求め、初期値の定理と最終値の定理で、はじめの値と最後の値を確かめましょう。
\bkjans{$F=\frac{10}{s}-\frac{10}{s+0.5}=\frac{5}{s(s+0.5)}$。$sF=\frac{5}{s+0.5}$。$s\to\infty$ で0、$s\to0$ で10。はじめ0\,V、最後10\,V。}
\end{jibun}

%% ---------------- まちがえやすいところ ----------------
\bksection{tsumazuki}{まちがえやすいところをチェック}

\begin{tsumazuki}{① $s$ のずらしの向き}
\bkng{$\mathcal L[e^{-t}\sin2t]=\dfrac{2}{(s-1)^2+4}$}
\bkok{$e^{-at}$ は $s+a$。$\mathcal L[e^{-t}\sin2t]=\dfrac{2}{(s+1)^2+4}$}
$e^{-t}e^{-st}=e^{-(s+1)t}$ なので、$s$ が $s+1$ になります。$e^{at}$（ふえる）なら $s-a$ です。
\end{tsumazuki}

\begin{tsumazuki}{② 名前を、そのままかけ算と読む}
\bkng{$\mathcal L[t\sin2t]=\mathcal L[t]\times\mathcal L[\sin2t]$}
\bkok{$t$ をかける行で、$-\dfrac{\dd F}{\dd s}$ に直す}
\[ \kixgNt{\mathcal L[t\sin2t]}=\kixgSt{2}{-\frac{\dd}{\dd s}\Bigl(\frac{2}{s^2+4}\Bigr)}=\frac{4s}{(s^2+4)^2} \]
$\mathcal L$ は「かける数」ではなく操作の名前です。かけ算の形は、性質の表（C・D・G・I）で直します。
\end{tsumazuki}

\begingroup\tcbset{equal height group=kixgEHT}%
\noindent\begin{minipage}[t]{0.485\linewidth}\vspace{0pt}
\begin{tsumazuki}{③ $H(t-t_0)$ を忘れる}
\raggedright\bkng{$\mathcal L[e^{-(t-1)}]=\dfrac{e^{-s}}{s+1}$}
\bkok{$e^{-(t-1)}H(t-1)$ のときだけ $\dfrac{e^{-s}}{s+1}$}
$H$ がないと、$t<1$ でも式が生きています。$e^{-(t-1)}=e\cdot e^{-t}$ で、変換は $\frac{e}{s+1}$ です。
\end{tsumazuki}
\end{minipage}\hfill
\begin{minipage}[t]{0.485\linewidth}\vspace{0pt}
\begin{tsumazuki}{④ $\delta$ を「高さ1」と思う}
\raggedright\bkng{$\delta(t)$ は、時刻0で高さ1}
\bkok{面積（力積）が1。高さは限りなく大きい}
高さ1で幅0なら、面積は0で、たたいたことになりません。
\end{tsumazuki}
\end{minipage}\par
\endgroup

\begingroup\tcbset{equal height group=kixgEHU}%
\noindent\begin{minipage}[t]{0.485\linewidth}\vspace{0pt}
\begin{tsumazuki}{⑤ 落ち着かない式に最終値の定理}
\raggedright\bkng{$\mathcal L[\sin t]=\dfrac{1}{s^2+1}$ なので、最後は $\lim\limits_{s\to0}\dfrac{s}{s^2+1}=0$}
\bkok{分母が0の点 $\pm j$ は左半分にないので使えない}
$\sin t$ は落ち着きません。
\end{tsumazuki}
\end{minipage}\hfill
\begin{minipage}[t]{0.485\linewidth}\vspace{0pt}
\begin{tsumazuki}{⑥ たたみこみを、ふつうのかけ算と思う}
\raggedright\bkng{$e^{-t}*e^{-2t}=e^{-3t}$}
\bkok{$e^{-t}*e^{-2t}=e^{-t}-e^{-2t}$}
かけ算になるのは、$s$ の世界でだけです。
\end{tsumazuki}
\end{minipage}\par
\endgroup

\begin{tsumazuki}{⑦ 分子の数を合わせずに表を引く}
\bkng{$\dfrac{1}{s^2+9}$ を見て、$\sin3t$}
\bkok{$\dfrac{1}{s^2+9}=\dfrac13\times\dfrac{3}{s^2+9}$ なので、$\dfrac13\sin3t$}
表の行 $\frac{\omega}{s^2+\omega^2}$ の分子は $\omega$ です。$\frac{1}{s+2}$ と $\frac{2}{s+2}$ も、$e^{-2t}$ と $2e^{-2t}$ でちがいます。
\end{tsumazuki}

\begin{tsumazuki}{⑧ 2つのずらしを取りちがえる}
\bkng{$e^{-t}f(t)$ の変換で、$F(s)$ に $e^{-s}$ をかける}
\bkok{$e^{-at}$ をかける → $s$ を $s+a$ に（$s$ のずらし）。$t_0$ おくれる → $e^{-st_0}$ をかける（時間のずらし）}
時間の世界で $e^{-at}$ をかけると $s$ の世界でずれ、時間の世界でずらすと $s$ の世界で $e^{-st_0}$ がかかります。
\end{tsumazuki}




\begin{column}{デルタ関数と、ヘビサイド・ディラック・シュワルツ}
短い時間にはたらく「たたく力」を、ふつうの関数のように式に入れる考えは、ヘビサイドが電気回路の計算で使っていました（9-5 のコラム）。物理学者ディラックは、量子力学の研究で、1927年の論文と1930年の本でこの関数を使い、$\delta$ という記号を広めました。値で決まらないこの「関数」は、数学者からは長く疑われましたが、1945年ごろにフランスの数学者シュワルツが、超関数（分布）という考えで正しく整えました。工学者の便利な道具が、あとから数学として認められた例の1つです。
\end{column}

\begin{motto}[5]
\kixgbox{9-4 のフーリエ変換の性質の表と、くらべてみる}
9-5 で、フーリエ変換は、ラプラス変換の $s$ を $j\omega$ にした所だと見ました（$t<0$ で0で、その積分が有限になるとき）。だから、この章の性質の表は、9-4 の性質の表と同じ形をしています。
\begin{bkcenter}
\begin{bktable}{colspec={Q[l,m,32mm] Q[c,m,42mm] Q[c,m,36mm]}, row{2-Z}={font=\small}, rowsep=1pt, colsep=3pt}
 & ラプラス変換（9-6） & フーリエ変換（9-4） \\
$e^{-at}H(t)$ & $\dfrac{1}{s+a}$ & $\dfrac{1}{a+j\omega}$ \\
微分 $f'$ & $s$ をかける（$-f(0)$ も） & $j\omega$ をかける \\
$t_0$ おくれる & $e^{-st_0}$ をかける & $e^{-j\omega t_0}$ をかける \\
たたみこみ $f*g$ & $F(s)\,G(s)$ & $F(\omega)\,G(\omega)$ \\
\end{bktable}
\end{bkcenter}
ちがいは、ラプラス変換の微分の行に $-f(0)$ があることです。ラプラス変換は $t=0$ から先だけを見るので、はじめの値が残ります。フーリエ変換は $-\infty$ から見るので、はじめの値という考えがありません。だから、はじめの値のある微分方程式を解くときは、ラプラス変換を使います。たとえば 9-4 で求めた $e^{-2t}H(t)$ のフーリエ変換 $\frac{1}{2+j\omega}$ は、この章の表の $\frac{1}{s+2}$ に $s=j\omega$ を入れたものと同じです。$\omega=1$ なら、どちらも $\frac{1}{2+j}=0.4-0.2j$ です。
\end{motto}

%% ---------------- 物理・工学ではこう使う ----------------
\begin{butsuri}
\kixgbox{ハンマーでたたいて、機械の性質を調べる}
\noindent\begin{minipage}[c]{0.64\linewidth}
\raggedright
機械の部品や建物のはりを、力をはかれるハンマー（インパルスハンマー）でたたき、ゆれをセンサーではかる試験があります。たたく力はほぼ $\delta(t)$ なので、はかったゆれはインパルス応答です。インパルス応答がわかれば、たたみこみで、どんな力を加えたときのゆれも計算できます。車や飛行機の部品の設計で使われています（物理 波-10）。
\end{minipage}\hfill
\begin{minipage}[c]{0.32\linewidth}
\centering
\begin{tikzpicture}[x=1mm, y=1mm]
  \fill[bkGray!40] (0,0) rectangle (36,3);
  \draw[bkInk, line width=0.7pt] (0,0) rectangle (36,3);
  \fill[bkInk] (3,-3) -- (6,-3) -- (4.5,0) -- cycle;
  \fill[bkInk] (30,-3) -- (33,-3) -- (31.5,0) -- cycle;
  \draw[bkInk, line width=1.4pt] (14,12) -- (24,20);
  \filldraw[fill=bkOrange!60, draw=bkInk, line width=0.7pt, rotate around={38:(14,12)}] (11,9) rectangle (17,14);
  \draw[bkRed, line width=1pt, -{Stealth[length=1.6mm]}] (13,7.5) -- (13,4);
  \filldraw[fill=bkBlue!30, draw=bkInk, line width=0.6pt] (22,3) rectangle (26,6);
  \node[font=\footnotesize, anchor=west] at (26.5,7) {センサー};
\end{tikzpicture}
\end{minipage}
\par\vspace{4pt}
\kixgbox{ホールの響きを、手をたたいて調べる}
コンサートホールで、短く大きな音（風船を割る音・ピストルの音）を出し、客席のマイクで響きを録音すると、ホールのインパルス応答になります。これとのたたみこみで、録音室で演奏した音を「そのホールで聞いたような音」にできます。音楽ソフトの「残響」の効果は、このしくみを使ったものがあります。
\par\vspace{4pt}
\kixgbox{シャワーと工場のむだ時間}
シャワーの温度調節がむずかしいのは、むだ時間 $e^{-sL}$ があるからです。温度を見て調節しても、効き目が $L$ 秒おくれて出るので、調節しすぎてしまいます。工場の長い配管やベルトコンベアの制御でも、むだ時間が設計のむずかしさになります（9-12）。
\end{butsuri}



%% ---------------- 章末問題 ----------------
\bkendproblems
\bkproblemgroup{基本}{5}

\bkproblem 次のラプラス変換を、表を引いて求めましょう\mbox{（\bkreidairef{1}）}。\par
（1）\nobreak $\mathcal L[t^4]$\par
（2）\nobreak $\mathcal L[\sin4t]$\par
（3）\nobreak $\mathcal L[\cos t]$
\bkpans{（1）\nobreak $\frac{24}{s^5}$。（2）\nobreak $\frac{4}{s^2+16}$。（3）\nobreak $\frac{s}{s^2+1}$。}

\bkproblem $\mathcal L[3-2t^2+5\sin t]$ を求めましょう\mbox{（\bkreidairef{2}）}。
\bkpans{$\frac{3}{s}-\frac{4}{s^3}+\frac{5}{s^2+1}$。}

\bkproblem 次のラプラス変換を求めましょう\mbox{（\bkreidairef{3}）}。\par
（1）\nobreak $\mathcal L[e^{-3t}\sin t]$\par
（2）\nobreak $\mathcal L[e^{-3t}\cos t]$\par
（3）\nobreak $\mathcal L[e^{-2t}\sin4t]$
\bkpans{（1）\nobreak $\frac{1}{(s+3)^2+1}=\frac{1}{s^2+6s+10}$。（2）\nobreak $\frac{s+3}{s^2+6s+10}$。（3）\nobreak $\frac{4}{(s+2)^2+16}=\frac{4}{s^2+4s+20}$。}

\bkproblem 次の関数のラプラス変換を求めましょう\mbox{（\bkreidairef{4}）}。\par
（1）\nobreak 3秒後にスイッチを入れて 12\,V をつなぐ電圧 $12H(t-3)$\par
（2）\nobreak $2H(t-5)$\par
（3）\nobreak $3\delta(t-2)$
\bkpans{（1）\nobreak $\frac{12e^{-3s}}{s}$。（2）\nobreak $\frac{2e^{-5s}}{s}$。（3）\nobreak $3e^{-2s}$。}

\bkproblem $\ddot x+6\dot x+10x=0$、はじめ $x(0)=2$、$\dot x(0)=0$ の両辺をラプラス変換して、$X(s)$ を求めましょう\mbox{（\bkreidairef{5}）}。
\bkpans{$(s^2X-2s)+6(sX-2)+10X=0$。$X(s)=\frac{2s+12}{s^2+6s+10}$。}

\bkproblem $x'+3x=2\delta(t)$、はじめ $x=0$ の両辺を変換して $X(s)$ を求め、表を逆に引いて $x$ を求めましょう\mbox{（\bkreidairef{6}）}。
\bkpans{$sX+3X=2$。$X=\frac{2}{s+3}$。$x=2e^{-3t}$。}

\bkproblem $f(t)=5-5e^{-4t}$ を変換し、初期値の定理と最終値の定理で、はじめの値と最後の値を確かめましょう\mbox{（\bkreidairef{11}）}。
\bkpans{$F=\frac{5}{s}-\frac{5}{s+4}=\frac{20}{s(s+4)}$。$sF=\frac{20}{s+4}$。$s\to\infty$ で0、$s\to0$ で5。}

\bkproblemgroup{チャレンジ}{5}

\bkproblem 次の電圧のラプラス変換を求めましょう\mbox{（\bkreidairef{7}）}。\par
（1）\nobreak はじめの1秒間だけ 4\,V をかけ、そのあと 0\,V にする電圧\par
（2）\nobreak 時刻2秒から5秒までの間だけ 1\,V をかけ、ほかの時刻は 0\,V の電圧
\bkpans{（1）\nobreak $4H(t)-4H(t-1)$。$\frac{4(1-e^{-s})}{s}$。（2）\nobreak $H(t-2)-H(t-5)$。$\frac{e^{-2s}-e^{-5s}}{s}$。}

\bkproblem $\mathcal L[t\sin3t]$ と $\mathcal L[t\sin4t]$ を求めましょう\mbox{（\bkreidairef{8}）}。
\bkpans{$t\sin3t$：$F=\frac{3}{s^2+9}$、$-\frac{\dd F}{\dd s}=\frac{6s}{(s^2+9)^2}$。$t\sin4t$：$F=\frac{4}{s^2+16}$、$-\frac{\dd F}{\dd s}=\frac{8s}{(s^2+16)^2}$。}

\bkproblem $e^{-t}*e^{-3t}$ を、たたみこみの積分と、$s$ の世界のかけ算で求め、同じになることを確かめましょう\mbox{（\bkreidairef{10}）}。
\bkpans{積分：$e^{-t}\int_0^te^{-2t_1}\,\dd t_1=\frac{e^{-t}-e^{-3t}}{2}$。かけ算：$\frac{1}{(s+1)(s+3)}$。$\frac12(\frac{1}{s+1}-\frac{1}{s+3})=\frac{1}{(s+1)(s+3)}$。同じ。}

\bkproblem $f(t)=\dfrac{1}{1+t}$ のラプラス変換の、$s=2$ のときの値を、表計算の台形（$\Delta t=0.1$、$t=10$ まで）で求めましょう\mbox{（\bkreidairef{9}）}。
\bkpans{かけた値は $\frac{1}{1+t}e^{-2t}$。$t=0$ で1、0.1 で0.7443…。台形の和は約0.364（$\Delta t=0.01$ で0.3614）。}

\begin{furikaeri}
\begin{checklist}
\item 表計算で $f(t)e^{-st}$ を台形でたし、変換表の値と同じになることを確かめられる。\bkfurref{bk:9-6:menseki}
\item $t^n$・$\sin$・$\cos$ と、$e^{-at}$ がかかった形の変換表の行を、部分積分・オイラーの公式・$s$ のずらしで作れる。\bkfurref{bk:9-6:tn}
\item 場合分けの表9-6と流れ図を使って、おくれ・パルス・デルタ関数・たたみこみをふくむ関数を変換し、初期値・最終値の定理で確かめられる。\bkfurref{bk:9-6:baai}
\end{checklist}
\end{furikaeri}

\bkfukidashi{mina}{変換表がそろった！　次は、部分分数に分けて、表を逆に引き、微分方程式の答えを出そう。}
